% \documentclass[journal]{IEEEtran}
\documentclass[letterpaper, 10 pt, conference]{ieeeconf}  % Comment this line out
                                                          % if you need a4paper
%\documentclass[a4paper, 10pt, conference]{ieeeconf}      % Use this line for a4
                                                          % paper

\IEEEoverridecommandlockouts                              % This command is only
                                                          % needed if you want to
                                                          % use the \thanks command
\overrideIEEEmargins

% *** MATH PACKAGES ***
%
\usepackage{amsfonts}
\usepackage{amsmath}
\usepackage{amssymb}
\newtheorem{theorem}{\textbf{Theorem}}
\newtheorem{corollary}{\textbf{Corollary}}
\newtheorem{lemma}{\textbf{Lemma}}
\newtheorem{proposition}{\textbf{Proposition}}
\newtheorem{definition}{\textbf{Definition}}
\newtheorem{assum}{\textbf{Assumption}}
\newtheorem{remark}{\textbf{Remark}}
% \usepackage{algorithmic}
% \usepackage[linesnumbered,ruled,vlined]{algorithm2e} 
\usepackage{makecell}
\usepackage{hyperref}
\usepackage{algorithm}
\usepackage{algpseudocode}

\usepackage{cellspace} 
\setlength{\cellspacetoplimit}{3pt}
\setlength{\cellspacebottomlimit}{3pt}
\usepackage{graphicx}
\usepackage[table]{xcolor}





\begin{document}
%
% paper title
% Titles are generally capitalized except for words such as a, an, and, as,
% at, but, by, for, in, nor, of, on, or, the, to and up, which are usually
% not capitalized unless they are the first or last word of the title.
% Linebreaks \\ can be used within to get better formatting as desired.
% Do not put math or special symbols in the title.
\title{\color{red}Feasibility Analysis and Constraint Selection \\ in Optimization-Based Controllers}


\author{ Panagiotis~Rousseas$^{1}$,~\IEEEmembership{Member,~IEEE,}  Haejoon~Lee$^{2}$,~\IEEEmembership{Graduate Student Member,~IEEE,} \\ Dimos V. Dimarogonas$^{1}$,~\IEEEmembership{Fellow,~IEEE} and Dimitra Panagou$^{2,3}$,~\IEEEmembership{Senior Member,~IEEE} 
\thanks{Haejoon Lee and Dimitra Panagou would like to acknowledge the support by the National Science Foundation (NSF) under Award Number 1942907 and the Air Force Office of Scientific Research (AFOSR) under Award No. FA9550-23-1-0163.}% <-this % stops a space
% \thanks{$^\dagger$Both authors have equal contribution.}
\thanks{$^{1}$Division of Decision and Control Systems,
    School of Electrical Engineering and Computer Science,
        KTH Royal Institute of Technology, Stockholm, Sweden
        {\tt\small dimos@kth.se, rousseas@kth.se}}%
\thanks{$^{2}$Department of Robotics,
        University of Michigan, Ann Arbor, MI, USA
        {\tt\small haejoonl@umich.edu}}%
\thanks{$^{3}$Department of Aerospace Engineering,
        University of Michigan, Ann Arbor, MI, USA
        {\tt\small dpanagou@umich.edu }}%
}


% make the title area
\maketitle

% As a general rule, do not put math, special symbols or citations
% in the abstract or keywords.
\begin{abstract}
\color{red}
Control synthesis under constraints is at the forefront of research on autonomous systems, in part due to its broad application from low-level control to high-level planning. Finding controllers is typically cast as a constrained optimization problem. Assessing feasibility of the constraints and selecting among subsets of feasible constraints is a challenging yet crucial problem. In this work, we provide a novel theoretical analysis that yields necessary and sufficient conditions for feasibility assessment. Concurrently, we develop theory-based methods for online feasible constraint selection. Through a series of simulations, we demonstrate that our algorithms achieve performance comparable to state-of-the-art methods while offering improved computational efficiency. Importantly, our theoretical results pave the way for formally addressing various problems relating to constrained control. 
% In this work, we propose a heuristic-based method that resolves infeasibility at every time step by selectively disregarding a subset of soft constraints based on the past values of the Lagrange multipliers. Compared to existing approaches, our method requires the solution of a smaller optimization problem to determine feasibility, resulting in significantly faster computation. Through a series of simulations, we demonstrate that our algorithm achieves performance comparable to state-of-the-art methods while offering improved computational efficiency.
\end{abstract}

% Note that keywords are not normally used for peerreview papers.
% \begin{IEEEkeywords}
% IEEE, IEEEtran, journal, \LaTeX, paper, template.
% \end{IEEEkeywords}






% For peer review papers, you can put extra information on the cover
% page as needed:
% \ifCLASSOPTIONpeerreview
% \begin{center} \bfseries EDICS Category: 3-BBND \end{center}
% \fi
%
% For peerreview papers, this IEEEtran command inserts a page break and
% creates the second title. It will be ignored for other modes.
\IEEEpeerreviewmaketitle



\section{Introduction}

% Many autonomous systems are deployed to accomplish tasks under numerous constraints arising from safety, performance, and operational requirements. Optimization-based or constrained controllers have recently received significant attention for systematically enforcing these requirements by computing control actions that satisfy constraints while optimizing performance objectives. However, constraints may sometimes be mutually incompatible, rendering the optimization problem infeasible and putting the system at risk of failure. Therefore, providing efficient and principled mechanisms for handling such incompatible constraints online is essential for the safe and trustworthy deployment of autonomous systems.


Enabling recursive feasibility of optimization-based controllers, and principled relaxation of constraints, as/if needed, is an essential capability towards trustworthy human-machine teaming, in the sense that assessing feasibility and relaxing constraints if/as needed in an explainable manner will enable humans and machines to make informed, trusted decisions, and collaborate efficiently. Hence, ensuring the safety, robustness and effectiveness of human-in-the-loop systems requires the development of novel methodologies that assess the feasibility of the underlying optimization problems, and that guide the relaxation of the constraints if the optimization problem is deemed infeasible.

The problem of constrained optimization has been extensively studied in the literature \cite{book_optim}. Finding the maximal subset of feasible constraints \cite{chinneck} is known to be an NP-hard problem; hence approximate solutions and heuristics are usually employed \cite{782859,Chinneck2010-qg}. A standard method to find feasible subsets is by adding relaxation variables (slack) to the constraints, accompanied by a suitable modification of the cost function \cite{Chinneck1996}. However, this results in a significant increase in the problem's size in the presence of multiple constraints, and in general does not guarantee that the maximal feasible set is found. Towards modeling the feasibility of multiple constraints in the presence of priority specifications, the authors in \cite{Dubois1996} propose a solution based on possibility theory built on search-based methods. 
\par
When it comes to path/trajectory planning and control of dynamical systems, a variety of methodologies aim to handle constraints. Techniques such as Model Predictive Control (MPC) \cite{Schwenzer2021}, Reference Governors \cite{GARONE2017306} and Control Barrier Functions (CBFs) \cite{ames2016control} enforce the satisfaction of constraints during the generation of control input trajectories, either over finite horizons or pointwise in time and state; however, guaranteeing recursive feasibility remains a challenge. More recently, heuristics for constraint selection based on Lagrange multipliers over dynamical system trajectories induced by CBF-QPs have been proposed \cite{hardik2024}. For high-level path planning under constraints, or more generally specifications (encoded e.g., via temporal logics), the case of infeasible/incompatible specifications and how to minimally relax or violate them has also been considered \cite{Rahmani2020WhatTD, 7989177, Wongpiromsarn2020MinimumViolationPF, 6579837,10342094, 10.1145/3450267.3450542}. However, such techniques typically neglect constraints due to dynamics, sensing and actuation limitations.
\par 
{\color{blue}More relevant to this work, in \cite{Molnar,cohen2025compatibilitymultiplecontrolbarrier} the authors investigate conditions for CBF-based constraints to be collectively feasible. Their approach can be used to guide control synthesis, with a focus on box input constraints (i.e., the $m$-dimensional input is constrained to lie within an $m$-parallelepiped). In contrast, we propose a condition that can be evaluated numerically and online to assess feasibility of any constraints that are linear w.r.t. the input. Our condition differs as it boils down to the solution of a LP, whereas the condition in \cite{Molnar} should hold over infinite many points. We will go into further detail on the differences between the former methods and ours' in the sequel.}
\par 
Notably, finding feasible control input trajectories under constraints typically requires computationally-efficient methods, i.e., methods that can find a solution fast enough for online implementation. Therefore, while many modern optimization programming tools do provide feasibility analysis, these may not scale with the size of the problem. Furthermore, in certain applications it might be desirable to avoid spending computational resources on the entire large-scale optimization problem to determine feasibility of its current instantiation. In our recent work \cite{rousseas2025feasibilityevaluationquadraticprograms} we developed a methodology for assessing the feasibility of Quadratic Programs (QPs), which is a standard form in constrained control synthesis such as in Model Predictive Control (MPC) and control synthesis subject to Control Barrier Functions (CBFs). Our approach casts the feasibility determination of the QP into a simpler Linear Program (LP), which can be solved and assessed more efficiently than the original QP. The proposed approach can then be used as a method for evaluating scenarios and determining in real-time which ones will be feasible over a horizon in the future. 
\par 
{
\color{blue}
Nevertheless, several problems persist. More specifically, \cite{rousseas2025feasibilityevaluationquadraticprograms} focused only on QP-based controllers by leveraging Lagrange Multipliers (LMs). At the same time, more general existing approaches \cite{hardik2024} also focus on LMs. The main technical issue with such treatments is that LMs capture the interplay between the optimal solution and the constraints' set. Analyzing (in)feasibility should be based only on the structure of the constraints, irrespective of the optimal solution to be extracted. While there exist methods that only consider the aforementioned structure \cite{Chinneck2010-qg,chinneck}, these involve constraint relaxation via auxiliary variables; the magnitudes of the former model the degree of constraint satisfaction or violation. However, relaxing the constraints fails to decouple two critical aspects: 1) how many and which constraints are enforced, versus 2) the degree of satisfaction or violation. More precisely, such algorithms may result in solutions that enforce few constraints (i.e., minimal violation of a large number of constraints), rather than enforcing as many constraints as possible \cite{chinneck_textbook}.    
}

\textit{\textbf{Contribution}}:
{\color{blue}
In view of the above, and building upon our earlier work on feasibility checking for QPs, the contribution of this work is twofold: 
\begin{enumerate}
    \item First, in contrast to existing approaches based on relaxation of the constraint set's inequalities, we provide a novel analysis for extracting sufficient and necessary conditions for feasibility (Sec. \ref{sec:method}), and
    \item Second, we employ the aforementioned analysis in order to provide online constraint selection algorithms for constrained control of dynamical systems (Sec. \ref{sec:conf_selec_methods}). 
\end{enumerate}
While the theoretical analysis of the sequel is employed for online constraint selection, its importance goes beyond this problem; we believe that our necessary and sufficient conditions can be further used for tackling the maximum feasible set problem (maxFS) \cite{chinneck}, feasible space monitoring \cite{hardik2024}, and for rendering infeasible problems feasible through control-based approaches. These directions will be explored in the future based on the herein foundational analysis.     
% Building upon our earlier work on feasibility checking for QPs, in this paper we provide: 1) A more general approach (based on Farkas' Lemma \cite{Gale1951}) that concentrates on the constraint set rather than focusing on the specific QP structure, 2) An analysis that results in method for evaluating feasibility of a set of constraints numerically through a linear program, 3) Theory-inspired methods for online constraint selection based on properties of the constraint set. Our method is demonstrated to be numerically efficient for online applications where constraint feasibility may be crucial to ensure system safety and satisfactory performance.
}
\par 
{\color{red}
\textit{\textbf{Organization}}: In Section~\ref{sec:prob}, we present some preliminary notions and formulate our problem. In Sec.~\ref{sec:method}, the theoretical elements of our approach are outlined, followed by Sec.~\ref{sec:conf_selec_methods} where two algorithms for online feasible constraint selection are presented. Finally, we provide simulations of the proposed algorithms in Sec.~\ref{sec:sims} with comparisons to two existing methodologies. Finally, some concluding remarks and future research directions are provided in Sec. \ref{sec:concl}. 
}
\par 
{\color{blue}
\textit{\textbf{Notation}}:
Let $\mathbb{N}$ denote the set of natural numbers, i.e., $\mathbb{N} = \{0, 1, 2, \dots\}$, $\mathbb{R}_{\geq 0}, \mathbb{R}_{< 0}$ denote a set of non-negative and strictly negative real numbers respectively. Similarly, for $n\in\mathbb{N}$, let $\mathbb{R}^n_{\geq 0}, \mathbb{R}^n_{< 0}$ denote $n$-dimensional vectors with non-negative and strictly negative real elements respectively.
Let $A\in{\color{blue}\mathbb{R}^{m\times c}}$ and $B\in\mathbb{R}{\color{blue}^c}$ denote a real-valued matrix and vector respectively. Then, $A_i\in\mathbb{R}^{m}$ denotes the $i$-th column of matrix $A$, and $B_i \in \mathbb{R}$ denotes the $i$-th element of vector $B$. Given two vectors $x \in \mathbb{R}^n,y\in\mathbb{R}^m$ let $(x,y) \in \mathbb{R}^{n+m}$ denote the stacked vector $(x,y) \triangleq [x^\top,y^\top]^\top$. Let $I_m$ denote the $m \times m$ identity matrix. We use ${1}_m$ and ${0}_m$ to denote $m$-dimensional vectors of ones and zeros, respectively. Let $\{0,1\}{\color{blue}^c}, c\in\mathbb{N}$ denote the set of vectors with $c$ elements consisting of either ones or zeros. 
}

\section{Problem Formulation}
\label{sec:prob}
\begin{figure}[h!]
    \centering
    \includegraphics[trim={2.0cm 4.cm 2.0cm 4cm},clip,width=1\linewidth]{figures/theory3.1.pdf}
    \caption{Overview of the Problem and its connection to the methodology (Secs. \ref{sec:method}, \ref{sec:conf_selec_methods}).}
    \label{fig:prob}
\end{figure}
\subsection{Preliminaries}
Consider the dynamical system:
\begin{equation}\label{eq:dynamics}
    \dot{x} = f(x) + g(x)u,
\end{equation}
where $x\in\mathbb{R}^n$ denotes the state vector, $f:\mathbb{R}^n\rightarrow\mathbb{R}^n,\ g:\mathbb{R}^n\rightarrow\mathbb{R}^{n\times m}$ are locally Lipschitz continuous functions and $u\in\mathbb{R}^m$ is the control input vector. Furthermore, for some time instance $t\in\mathbb{R}_{\geq0}$ consider {\color{blue}$c\in\mathbb{N}\setminus\{0\}$} affine constraints w.r.t. the input $u\in\mathbb{R}^m$ given by:
\begin{equation}\label{eq:affine_cnstr}
   %\mathcal{U} \triangleq \left\{u\in\mathbb{R}^m | A^\top(x,t)u \leq B(x,t) \right\},\\
A_i^\top(x,t)u \leq B_i(x,t), \quad i\in \mathcal C\triangleq \lbrace 1,\dots,c\rbrace,
\end{equation}
where $A_i : \mathbb{R}^n \times \mathbb{R}_{\geq 0}\rightarrow \mathbb{R}^m$ and $B_i: \mathbb{R}^n \times \mathbb{R}_{\geq 0}\rightarrow\mathbb{R}$ are Lipschitz continuous wrt to $x$ and $t$. %$A:\mathbb{R}^n\times \mathbb{R}_{\geq0}\rightarrow{\color{blue}{\color{blue}\mathbb{R}^{m\times c}}},\ B:\mathbb{R}^n\times \mathbb{R}_{\geq0}\rightarrow{\color{blue}\mathbb{R}^{c}}$ are locally Lipschitz continuous with respect to $x$ and $t$.
  %The constraints encoded {\color{blue}in $A,B$} consist of two subsets, namely hard constraints, which should always be satisfied, and soft constraints, which can be disregarded in case of infeasibility. 
  Without loss of generality, we assume that $n_h$ of the constraints are hard and $n_s$ of the constraints are soft, so that $n_h+n_s=c$, $\mathcal{C}_h=\{1,\dots,n_h\}$, $\mathcal C_s=\{n_h+1,\dots,c\}$. We assume:
\begin{assum}
\label{assum:hard_const_compatible}
    The set of hard constraints is always feasible, i.e., ${\color{blue}\forall }x\in\mathbb{R}^n, \forall t\geq 0$: 
    \begin{equation}
    \mathcal{U}_h \triangleq \left\{ u\in\mathbb{R}^m | \bigcap_{i\in\mathcal{C}_h} \{A_i^\top(x,t)u\leq B_i(x,t)\} \right\}\neq \emptyset,
    \end{equation}
    where $\mathcal{U}_h \subseteq \mathbb{R}^m$ denotes the admissible input set w.r.t. the hard constraints. %$A_i : \mathbb{R}^n \times \mathbb{R}_{\geq 0}\rightarrow \mathbb{R}^m$ denotes the $i$-th column of $A$ and $B_i: \mathbb{R}^n \times \mathbb{R}_{\geq 0}\rightarrow\mathbb{R}$ denotes the $i$-th element of $B$, where the arguments are dropped for brevity. 
\end{assum}



\subsection{Modeling Disregarded Constraints}

{
\par 
{\color{blue}
The soft constraints are divided into two sets, namely the set of enforced constraints and the set of disregarded constraints, indexed as $\mathcal{C}_e \subseteq \mathcal{C}_s$ and $\mathcal{C}_d \subseteq \mathcal{C}_s$ respectively, such that $\mathcal{C}_e \cap \mathcal{C}_d = \emptyset$ and $\mathcal{C}_e \cup \mathcal{C}_d = \mathcal{C}_s$. The disregarded constraints are defined as those are dropped from the constraint set \eqref{eq:affine_cnstr} to yield a feasible subset of constraints. In other words, only the enforced constraints are considered:
\begin{equation}
    \mathcal{U}_e \triangleq \left\{ u \in \mathbb{R}^m | \bigcap_{i \in \mathcal{C}_e} \{A_i^\top(x,t)u \leq B_i(x,t)\} \right\} .
\end{equation}
Hence the admissible input set $\mathcal{U}'\subset \mathbb{R}^m$ is $\mathcal{U}'\triangleq \mathcal{U}_h \cap \mathcal{U}_e$.
}
% Given a soft constraint, i.e., a pair consisting of a column of the matrix $A_i: \mathbb{R}^n \times \mathbb{R}_{\geq 0}\rightarrow\mathbb{R}^m$ and the corresponding element of the vector $B_i: \mathbb{R}^n \times \mathbb{R}_{\geq 0}\rightarrow\mathbb{R}$ where $i\in\mathcal{C}_s$ (the arguments are dropped for brevity), and a corresponding set of admissible inputs $\mathbb{R}^m \supseteq \mathcal{U}_s \triangleq \left\{ u\in\mathbb{R}^m | A_i^\top(x,t)u\leq B_i(x,t),\forall i\in\mathcal{C}_s \right\}$ a disregarded constraint is defined as follows:
% \begin{definition}\label{def:disreg_constr}
%     A constraint $A_i^\top(x,t) u \leq B_i(x,t),i\in\mathcal{C}_s$ is \textbf{disregarded} if it is dropped from the soft constraints' set, i.e.: $\mathbb{R}^m \supseteq \mathcal{U}_s'\triangleq  \left\{ u\in\mathbb{R}^m | A_i^\top(x,t)u\leq B_i(x,t),\forall i\in\mathcal{C}_e \right\}$, where $\mathcal{C}_e \triangleq \mathcal{C}_s \backslash \mathcal{C}_d$ denotes the index set of enforced constraints and $\mathcal{C}_d\subseteq \mathcal{C}_s$ denotes the index set of disregarded constraints. 
% \end{definition}
% Through the above definitions, note that $\mathcal{U} = \mathcal{U}_h \cap \mathcal{U}_s$ and denote the admissible input set, with the disregarded soft constraints indexed in $\mathcal{C}_d$ dropped as: $\mathcal{U}' \triangleq  \mathcal{U}_h \cap \mathcal{U}_s'$.



Let the $c$-dimensional \textbf{configuration} vector $P\in \mathcal{P}$, where $\mathcal P \triangleq \{0,1\}{\color{blue}^c}$, be such that the $i$-th element is $0$ when the $i$-th constraint is disregarded, and $1$ when the $i$-th constraint is enforced. Denote also 
\begin{subequations}
\begin{align}
\label{eq:AB}
A(x,t) &= \bigl [A_1(x,t) | \dots | A_c(x,t)\bigr] \in \mathbb R^{m\times c},\\ B(x,t)&=\begin{bmatrix} B_1(x,t) &\dots & B_c(x,t) \end{bmatrix}^\top \in \mathbb R^c,
\end{align}
\end{subequations}
where $A_i(x,t)$ and $B_i(x,t)$, $i\in \mathcal C$ are as in  \eqref{eq:affine_cnstr}.  
\begin{remark}
\color{blue}
    The elements of the configuration vector $P\in\mathcal{P}$ corresponding to hard constraints are always set equal to one, i.e., $P_i = 1, \forall i \in \mathcal{C}_h$.
\end{remark}
Then, the set of admissible inputs $\mathcal{U}'$ can be written as a family of sets:
\begin{equation}\label{eq:input_set:config}
    \mathcal{U}'(P) = 
    \left\{
        u \in \mathbb{R}^m 
        | 
        \textrm{diag}(P)
        A^\top(x,t)u \leq
        \textrm{diag}(P)
        B(x,t)
    \right\},
\end{equation}
where {\color{blue}$\textrm{diag}(P)\in \{0,1\}^{c\times c}$} denotes the matrix with $P\in\mathcal{P}$ in its diagonal and zeros elsewhere, which has the effect of nullifying the disregarded constraints indexed by $\mathcal{C}_s$ in $\mathcal{U}'$. 
{\color{blue}
For any configuration $P\in\mathcal{P}$, we define an index set $\mathcal{C} \supseteq \mathcal{C}_P \triangleq \{ i \in \mathcal{C} \;| \; P_i =1 \}$ that contains the indices of the enforced constraints in $\mathcal{U}'$. 

}
\par
The set of \textbf{feasible configurations} $\mathcal{P}_f$ is the set that contains all configurations $P\in\mathcal P$ whose input set is non-empty, i.e.: 
\begin{equation}\label{eq:feas_conf_set_def}
\begin{gathered}
    \mathcal{P} \supseteq \mathcal{P}_f \triangleq  
    \left\{ 
            P \in \mathcal{P} \; | \;  
            \mathcal{U}'(P)\neq \emptyset 
    \right\}.
\end{gathered}\end{equation}




% \begin{assum}\label{def:policy}
%     Given a {\color{blue} the constraint matrices $A,B:A\in{\color{blue}\mathbb{R}^{m\times c}},B\in\mathbb{R}{\color{blue}^c}$} and a corresponding \textbf{feasible} configuration vector $P\in\mathcal{P}_f$, there exists a mapping $\pi:\mathcal{P}_f\times{\color{blue}\mathbb{R}^{m\times c}}\times{\color{blue}\mathbb{R}^{c}}\rightarrow \mathbb{R}^m$ that is termed a \textbf{policy} and provides an input $u = \pi(P,A,B)$ to System \eqref{eq:dynamics}. The policy is also given in state-feedback form: {\color{blue}$\pi_{A,B}:\mathcal{P}_f\times\mathbb{R}^n\times\mathbb{R}_{\geq 0}\rightarrow\mathbb{R}^m$ and can be evaluated point-wise as: $\pi_{A,B}(P,x,t) \triangleq \pi(P,A(x,t),B(x,t))$.}  
% \end{assum}


% Given an admissible policy $\pi$, the trajectory of \eqref{eq:dynamics} from the initial state $\bar{x}\in\mathbb{R}^n$ under $\pi$, is denoted by $x_\pi:\mathbb{R}^n\times \mathbb{R}_{\geq0}\rightarrow\mathbb{R}^n$ and its value is given by $x_\pi = x_{\pi}(\bar{x},t)$. That is, $x_\pi$ is the solution to \eqref{eq:dynamics} for $u = \pi_{(A,B)}\left( P, x,t \right)$.
% Finally, we define the \textbf{level} of a configuration $L:\mathcal{P}\rightarrow\mathbb{N}$ as  $L(P) = \| P \|_1$. 

% \begin{remark}\label{rem:cbfqp}
%     An example of a class of admissible policies per Def.~\ref{def:policy} are CBFQP controllers, i.e.:
%     \begin{equation}\begin{gathered}
%        \pi_{(A,B)}(P,x,t)  = \underset{u\in\mathbb{R}^m}{\arg\min}\left\{ \|u - u_{ref}\|_2^2 \right\}\\
%         \textrm{s.t.} \quad \textrm{diag}(P)A(x,t)^\top u \leq\textrm{diag}(P)B(x,t),
%     \end{gathered}\end{equation}
% where $u_{ref}\in\mathbb{R}^m$ denotes a reference input to system~\eqref{eq:dynamics}.
% \end{remark}

{
\color{red}
\noindent\textbf{Problem Statement:} 
The goals of this work are: 1) Provide necessary and sufficient conditions for the (in)feasibility of collections of affine constraints \eqref{eq:affine_cnstr}, i.e., (non-)emptiness of the corresponding sets, and 2) Formulate methods for online constraint selection in the context of optimization-based controllers for dynamical systems (Fig. \ref{fig:prob}). 
}
}
\begin{figure*}
    \centering
    \includegraphics[trim={1.25cm 3.0cm 0.cm 3.5cm},clip,width=1\linewidth]{figures/theory3.pdf}
    \caption{Overview of the Theoretical Results in Sec. \ref{sec:method}.}
    \label{fig:theory}
\end{figure*}

\section{Methodology}\label{sec:method}
We now present our theoretical results for feasibility analysis, outlined in Fig. \ref{fig:theory}. 

\subsection{Analysis for Enforcing all of the Constraints}

In this Subsection, we analyze the case where all of the constraints encoded in \eqref{eq:affine_cnstr} are enforced. 
\subsubsection{\textbf{Feasibility Analysis \cite{rousseas2025feasibilityevaluationquadraticprograms}}}
{\color{blue}
In \cite{rousseas2025feasibilityevaluationquadraticprograms}, a Linear Program (LP) that assesses feasibility of a set of linear constraints \eqref{eq:affine_cnstr} was proposed:
}
\begin{equation}\label{eq:feas_lp:1}
    \begin{split}
        d^\star \in \underset{\lambda\in\Lambda}{\max}
        \left\{
            -B^\top \lambda 
        \right\},        
    \end{split}
\end{equation}
where $B := B(x,t)$, $A := A(x,t)$ are defined in \eqref{eq:AB}, 
$\Lambda = 
    \left\{ 
        \lambda \in \textrm{null}\left(A\right) \subset \mathbb{R}{\color{blue}^c} 
        |
        \lambda \geq 0
\right\}$. More specifically, the set $\mathcal U=\left\{u\in\mathbb{R}^m | A^\top u \leq B \right\}$ is non-empty iff $d^\star$ in \eqref{eq:feas_lp:1} is bounded. 
In \cite{rousseas2025feasibilityevaluationquadraticprograms}, {\color{blue}the proof was based on the assumption that the optimization cost is of  quadratic form.} Here this result is generalized to not be limited to quadratic programs (QPs) by demonstrating its connection to Farkas' Lemma \cite{Gale1951}:
\begin{lemma}[Farkas' Lemma \cite{Gale1951}]\label{thm:farkas}
    Given two matrices $A\in{\color{blue}\mathbb{R}^{m\times c}},B\in\mathbb{R}{\color{blue}^c}$, then exactly one of the following conditions hold. (1) $A^\top u \leq B$ has a solution $u\in\mathbb{R}^m$ (Condition 1). (2) $A \lambda =0 $ (i.e., $\lambda \in  \textrm{null}(A)$) has a solution $\lambda \geq 0$ with $B^\top \lambda <0$ (Condition 2).
\end{lemma}

We begin by showing a useful proposition:
\begin{proposition}\label{prop:lambda_zero}
    If a bounded {\color{blue} maximum} to \eqref{eq:feas_lp:1} exists, then {\color{blue}the maximizer is equal to} $\lambda^\star = [0,\cdots,0]^\top \in \mathbb{R}{\color{blue}^c}$.
\end{proposition}
\begin{proof}
    Assume for the sake of contradiction that $\exists \bar{\lambda}\neq [0,\cdots,0]^\top$ such that {\color{blue} $d^\star < \infty$ in} \eqref{eq:feas_lp:1} is maximum. Then, $\forall s \geq 0$, since $s A \bar{\lambda} = 0$ and $s\bar{\lambda} \geq 0$, it holds that  $s\bar{\lambda} \in\Lambda$. Define $d:\mathbb{R}_{\geq 0}\rightarrow \mathbb{R}$ given by {\color{blue}$d(s) \triangleq -s \cdot B^\top \bar{\lambda}$}. We now distinguish two cases: 1) If $-B^\top \bar{\lambda} > 0 $ then $\frac{\partial d}{\partial s} = -B^\top\bar{\lambda} >0 $ and hence $d=-sB^\top \bar{\lambda}$ grows unbounded as $s\rightarrow\infty$, contradicting boundedness of $d^\star$. 2) If  $-B^\top \bar{\lambda} \leq 0 $, then for $s^\star = 0 :  -s^\star B^\top \bar{\lambda}=0 \geq -B^\top \bar{\lambda}$. This shows that the maximizer is $s^\star \bar{\lambda} =  [0,\cdots,0]^\top $, contradicting the original assumption. We have thus established that if \eqref{eq:feas_lp:1} is bounded, its {\color{blue}maximizer} is $\bar{\lambda} = [0,\cdots,0]^\top \in \mathbb{R}{\color{blue}^c}$ with the {\color{blue}maximum equal to $d^\star = B^\top [0,\cdots,0]^\top =  0$}.
\end{proof}

\begin{theorem}\label{thm:feas_lp_farkas}
    Given two matrices $A\in{\color{blue}\mathbb{R}^{m\times c}},B\in\mathbb{R}{\color{blue}^c}$, the set $\mathcal{U} = \left\{u\in\mathbb{R}^m | A^\top u \leq B \right\}$ is non-empty iff \eqref{eq:feas_lp:1} is bounded. 
\end{theorem}
\begin{proof} 
    (Sufficiency) Assume that \eqref{eq:feas_lp:1} is bounded. This implies that Condition 2 in Farkas' Lemma \ref{thm:farkas} cannot hold. To see this, if $\exists \bar{\lambda} \geq 0 : A\bar{\lambda} =0 $ and $B^\top \bar{\lambda} < 0 \Rightarrow-B^\top \bar{\lambda} > 0$, then following the arguments of Prop. \ref{prop:lambda_zero}, this contradicts the boundedness of \eqref{eq:feas_lp:1}. Therefore, Condition 1 of Farkas' lemma must hold, implying that $\mathcal{U}$ is non-empty.
    \par 
    % We will show this by contraposition, i.e., we need to show that if \eqref{eq:feas_lp:1} is not bounded, then the set $\mathcal{U}$ is empty. 
    (Necessity) Assume that $\mathcal{U}$ is non-empty, which implies that Condition 1 of Farkas' lemma holds. Since Condition 2 in \ref{thm:farkas} cannot then hold, $\nexists \lambda \in \textrm{null}(A) : \lambda \geq 0$ and $-B^\top \lambda > 0$. Therefore, $\forall \lambda \in \textrm{null}(A), \lambda \geq 0$ it holds that $-B^\top \lambda \leq  0$. Then, following the arguments of Prop. \ref{prop:lambda_zero},  $\lambda = [0,\cdots,0]^\top \in \mathbb{R}{\color{blue}^c}$ is a {\color{blue}maximizer} of \eqref{eq:feas_lp:1} and the maximum is $-B^\top \lambda  =  0$. 
\end{proof}
\begin{remark}\label{rem:feas_lp:contra}
    By contraposition and Thm. \ref{thm:feas_lp_farkas}, if \eqref{eq:feas_lp:1} is unbounded, then the set $\mathcal{U}$ is necessarily empty.
\end{remark}
\textcolor{orange}{Thm.~\ref{thm:feas_lp_farkas} generalizes \cite[Thm.~1]{rousseas2025feasibilityevaluationquadraticprograms}, as the same result holds regardless of the structure of the cost function.} From Thm. \ref{thm:feas_lp_farkas}, the configuration $P\in\mathcal{P}$ is feasible if {\color{blue}$d^\star<\infty$ where}:
\begin{equation}\label{eq:feas_lp:2}
    \begin{split}
        d^\star \in \underset{\lambda\in\Lambda}{\max}
        \left\{
            -B^\top \textrm{diag}(P)\lambda 
        \right\},        
    \end{split}
\end{equation}
where 
$\Lambda = 
    \left\{ 
        \lambda \in \textrm{null}\left( \textrm{diag}(P)A\right) \subset \mathbb{R}{\color{blue}^c} 
        |
        \lambda \geq 0
\right\}$.
Otherwise, i.e., if {\color{blue}$d^\star$ in} \eqref{eq:feas_lp:2} is unbounded,  by Rem. \ref{rem:feas_lp:contra} $P$ is infeasible.


\subsubsection{\textbf{Feasibility Analysis using the Polar Cone}}
Here, we employ Thm.~\ref{thm:feas_lp_farkas} in order to form a criterion for assessing the impact of each constraint towards rendering the set $\mathcal{U}$ empty.
Denote $\mathbb{N} \ni {\color{blue}k} \triangleq \textrm{dim}\left( \textrm{ker} \left( A\right) \right)$, and let $\{n_1,\dots,n_c\}$ be a basis of $\textrm{null}(A)$, where $n_i \in \mathbb{R}^c, i \in \{1,\cdots,k\}$ are the basis vectors. Define $N_A = \bigl[ n_1 | n_2 | \cdots | n_k\bigr] \in {\color{blue}\mathbb{R}^{c\times k}}$, {\color{red}which we refer to as the nullspace basis}. Then there exists a unique $\mu \in {\color{blue}\mathbb{R}^k}$ such that $\textrm{null}(A)\ni\lambda = N_A \mu$, and 
%Note that by expressing $\lambda \in \textrm{null}(A)$ in terms of the nullspace basis, i.e., $\textrm{null}(A)\ni\lambda = N_A \mu$, for $\mu \in {\color{blue}\mathbb{R}^k}$, 
\eqref{eq:feas_lp:1} can be written equivalently as:
\begin{equation}\label{eq:feas_lp:star}
    \begin{split}
        d^\star \in &\underset{\mu\in{\color{blue}\mathbb{R}^k}}{\max}
        \left\{
            -B^\top N_A \mu 
        \right\},        \\
        &\textrm{s.t.: } N_A \mu\geq 0.
    \end{split}
    \tag{\ref{eq:feas_lp:1}*}
\end{equation}
By Thm. \ref{thm:feas_lp_farkas}, if {\color{blue}$d^\star$ in} \eqref{eq:feas_lp:star} is bounded (resp. unbounded), $\mathcal{U}$ is nonempty (resp. empty). The LP \eqref{eq:feas_lp:star} is a conic linear problem. Given the cone $\mathcal{N} \triangleq \{ \mu \in {\color{blue}\mathbb{R}^k} | N_A \mu \geq 0\}$, we consider the \textit{polar cone}, i.e.:
\begin{equation}\label{eq:polar_cone}
    \mathcal{N}_C = 
    \left\{ 
        \nu \in {\color{blue}\mathbb{R}^k} | \nu^\top\mu \leq 0, \forall \mu \in \mathcal{N}
    \right\}.
\end{equation}
We now show that for {\color{blue}$d^\star$ in} \eqref{eq:feas_lp:star} to be bounded, the vector $- N_A^\top B \in {\color{blue}\mathbb{R}^k}$ needs to lie within the polar cone $\mathcal{N}_C$. 
\begin{proposition}\label{prop:cone}
    The {\color{blue}maximum of the} LP \eqref{eq:feas_lp:star} is bounded iff $-B_A\in \mathcal{N}_C$, with $B_A \triangleq  N_A^\top  B$. 
\end{proposition}
\begin{proof}
    (Sufficiency) Assume that $-B_A\in \mathcal{N}_C$. By \eqref{eq:polar_cone}: 
    \begin{equation}
    \begin{split}
        -B^\top N_A \mu \leq 0, &\forall \mu \in \mathcal{N}\Leftrightarrow \\
         -B^\top N_A \mu \leq 0, &\forall \mu \in {\color{blue}\mathbb{R}^k}| N_A \mu \geq 0
    \end{split}
    \end{equation}
    Therefore, $\forall \mu \in {\color{blue}\mathbb{R}^k}| N_A \mu \geq 0$ it follows that $-B^\top N_A \mu \leq 0 \Rightarrow \max\left\{  -B^\top N_A \mu \right\}\leq 0$, and the {\color{blue}maximum of the} LP \eqref{eq:feas_lp:star} is bounded. 
    \par 
    (Necessity) We want to show that if $d^*<\infty$, then $-B_A \in \mathcal N_C $. We will prove this by contraposition, i.e., that if $-B_A \notin \mathcal{N}_A$, then {\color{blue}$d^*=\infty$}. For $-B_A \notin \mathcal{N}_A $, $\exists \bar{\mu}\in{\color{blue}\mathbb{R}^k}:N_A\bar{\mu}\geq 0$ such that $-B_A^\top \bar{\mu} > 0 \Leftrightarrow -B^\top  N_A \bar{\mu} > 0$.
    % \begin{equation}
    %     \begin{split}
    %         \exists \bar{\mu}\in{\color{blue}\mathbb{R}^k}:N_A\bar{\mu}\geq 0 \textrm{ such that:}  &-B_A^\top \bar{\mu} > 0 \Leftrightarrow -B^\top  N_A \bar{\mu} > 0.
    %     \end{split}
    %     \notag 
    % \end{equation}
    Now consider some $s\geq 0$ and define $d:\mathbb{R}_{\geq 0} \rightarrow \mathbb{R}$, given by $d(s) \triangleq  -s  B^\top N_A \bar{\mu} $. For $s\geq 0$ it holds that $N_A \bar{\mu} \geq 0 \Rightarrow N_A s \bar{\mu} = s N _A \bar{\mu}  \geq 0$. Therefore, the vector family $\mu = s\bar{\mu} \in {\color{blue}\mathbb{R}^k}$ satisfies the inequality in \eqref{eq:feas_lp:star} and $\frac{\partial d}{\partial s} =  -B^\top N_A \bar{\mu} > 0$. Therefore, $d$ grows unbounded as $s\rightarrow\infty$, implying that the {\color{blue}maximum of the} LP \eqref{eq:feas_lp:star} is unbounded, $d^\star=\infty$. Since  $-B_A \notin \mathcal{N}_C \Rightarrow d^\star = \infty$, this concludes the proof. 
\end{proof}

\subsubsection{\textbf{Feasibility Criterion}}
%In the previous subsection we have 
Proposition \ref{prop:cone} establishes that for \eqref{eq:feas_lp:star} to be feasible, the vector $B_A \triangleq -N_A^\top B $ must lie in $\mathcal{N}_C$. This is equivalent to feasibility of the following LP:
\begin{equation}\label{eq:polar_cone_B}
    \begin{gathered}
    \nu^\star \in \underset{\nu\in\mathbb{R}{\color{blue}^c}}{\arg\min}\left\{{1}^\top_c \nu \right\}, \\
    \textrm{s.t.: } N_A^\top \nu = B_A , \quad  \nu \geq 0,
    \end{gathered}    
\end{equation}
where ${1}_c = \begin{bmatrix}1 & 1 &\cdots&1\end{bmatrix}^\top\in\mathbb{R}{\color{blue}^c}$. The equality $ N_A^\top \nu = B_A$ in \eqref{eq:polar_cone_B} implies that the maximizer $\nu^\star$ (if it exists) is the projection of the gradient vector $B_A$ onto the column space of $\mathcal{N}_A$. We connect the LP \eqref{eq:polar_cone_B} to the LP \eqref{eq:feas_lp:star} through the following proposition:
\begin{proposition}\label{prop:lp_cone_feas}
    The maximum of the LP \eqref{eq:feas_lp:star} is bounded iff the LP \eqref{eq:polar_cone_B} is feasible. 
\end{proposition}
\begin{proof}
    Considering \eqref{eq:feas_lp:star}, its dual problem is:
\begin{equation}\label{eq:feas_lp:star:dual}
    \begin{gathered}
        \underset{\nu'\in\mathbb{R}{\color{blue}^c}}{\min}
        \left\{
            0 
        \right\},        \\
        \textrm{s.t.: } (-N_A)^\top \nu'= -B_A, \nu' \geq 0.
    \end{gathered}
\end{equation}
According to the weak duality principle, the maximizer of \eqref{eq:feas_lp:star} is bounded iff its dual \eqref{eq:feas_lp:star:dual} is feasible. Further, note that, for $\mathbb{R}{\color{blue}^c} \ni \nu \leq 0 \Rightarrow -{1}_c^\top\nu \geq 0 $ hence \eqref{eq:polar_cone_B} is bounded from below. Therefore, since the constraints in \eqref{eq:feas_lp:star:dual} and \eqref{eq:polar_cone_B} are the same, feasibility of the former implies feasibility of the later and vice-versa, concluding the proof.  
\end{proof}

{\color{blue}
\begin{remark}
    In \cite[Thm. 2]{Molnar}, \cite[Lemma 1]{cohen2025compatibilitymultiplecontrolbarrier}, the authors employ constraints of the form:
    \begin{equation}
        c_i(x) + d_i(x)^\top u \ge 0, \quad i \in \mathcal{I} \subset \mathbb{N},\notag 
    \end{equation}
    where $c_i : \mathbb{R}^n \to \mathbb{R}$ and $d_i : \mathbb{R}^n \to \mathbb{R}^m$. The constraints are compatible for
 (1) at
    $x \in \mathbb{R}^n$ if and only if \textbf{\textit{for all}} $\lambda_i \ge 0$ the following holds:
\begin{equation}\label{eq:molner_condition_feas}
        \sum_{i \in \mathcal{I}} \lambda_i d_i(x) = 0
        \quad \Longrightarrow \quad
        \sum_{i \in \mathcal{I}} \lambda_i c_i(x) \ge 0.
        % \notag
    \end{equation}
    Consider now a fixed state $x\in\mathbb{R}^n$. Comparing \eqref{eq:molner_condition_feas} (provided in \cite[Thm. 2]{Molnar}) to \eqref{eq:feas_lp:1}, $\sum_{i \in \mathcal{I}} \lambda_i d_i(x) = 0$ corresponds to $\lambda \in\textrm{null}(A) \Leftrightarrow \lambda = N_A \mu$ which together with the condition $\lambda_i \ge 0$ yield $\lambda \in\Lambda$, while $\lambda_i c_i(x)$ is exactly $B^\top \lambda$. Then, $\lambda_i c_i(x) \ge 0$ ensures that $\underset{\lambda \in\Lambda}{\max}\left\{-B^\top \lambda\right\}\leq 0 $ is bounded from above. Hence, the two conditions are equivalent. 
     However, note that even for a single $x\in\mathbb{R}^n$, the condition in \eqref{eq:molner_condition_feas} needs to be verified $\forall \lambda \geq 0$. In contrast, the proposed LP \eqref{eq:feas_lp:star} offers a direct way of evaluating feasibility. 
    % However, note that \eqref{eq:molner_condition_feas} should hold $\forall \lambda \geq 0$. Hence, our formulation covers \eqref{eq:molner_condition_feas}, with the main difference that the proposed LP \eqref{eq:feas_lp:star} offers a direct way of evaluating feasibility, while per Rem. \ref{rem:B_cone_components} the solution to the LP can also be used to assess whether a constraint renders the problem infeasible. 
\end{remark}
}
{\color{red}
\begin{remark}\label{rem:B_cone_components}
    % Each one of the components {\color{blue}of the solution to \eqref{eq:polar_cone_B}} $\nu^\star \in \mathbb{R}{\color{blue}^c}$ corresponds each of the $c$ row-vectors of the nullspace basis {\color{blue}in the columns of} $N_A\in{\color{blue}\mathbb{R}^{c\times k}}$. 
    By virtue of the equality constraint $N_A^\top \nu^\star = B_A$, each one of the components of the solution to \eqref{eq:polar_cone_B} $\nu^\star \in \mathbb{R}{\color{blue}^c}$ corresponds to each of the $c$ row-vectors of the matrix $N_A\in{\color{blue}\mathbb{R}^{c\times k}}$.
    \par 
    Note that due to the constraints in \eqref{eq:polar_cone_B}, $\nu^\star \geq 0$ element-wise. In contrast, if no such positive solution exists, (i.e., if any element of $\nu^\star$ should be negative to satisfy the equality constraint in \eqref{eq:polar_cone_B}), this implies emptiness of $\mathcal{U}$, and through Prop. \ref{prop:cone} it further implies that $B_A$ lies \textbf{outside} the polar cone \eqref{eq:polar_cone}. In case some elements of $\nu^\star $ are equal to zero, this implies that $B_A$ lies exactly on the boundary of the polar cone and more specifically on the plane(s) formed by $(N_A)_i^\top \nu = 0, \forall i\in \{1,\cdots,c\}: \nu^\star_i = 0$.
    \par 
    Therefore, we conclude that smaller magnitudes of the elements of $\nu^\star$, $|\nu^\star_i|, i\in \{1,\cdots,c\}$ indicate that $B_A$ is closer to the boundary of the polar cone, i.e., the distance $\underset{z \in \partial \mathcal{N}_C}{\min}\{\| z - B_A \| \}$ decreases, and the set $\mathcal{U}$ tends to become empty. 
    % The magnitude of each of the components of $\nu^\star$ then indicates how far the vector $B_A$ is from leaving the polar cone $\mathcal{N}_A$ and hence inducing emptiness of $\mathcal{U}$. 
\end{remark}
}
To formalize the above, we prove the following: 
{\color{blue}
\begin{proposition}\label{prop:minimal_sol}
        Consider the solution to \eqref{eq:polar_cone_B} denoted by $\nu^\star$ and the cone $\mathcal{N} \triangleq \{ \mu \in {\color{blue}\mathbb{R}^k} | N_A \mu \geq 0\}$. Assume that some of the rows of $N_A$ are redundant. Let $N_A^i \in \mathbb{R}^{1\times k}, i\in \mathcal{C}$ denote the rows of $N_A$ and also let $\mathcal{C}_r \subset \mathcal{C}$ denote the indices of the redundant rows of $N_A$ (redundancy means that: $N_A^i \mu \geq 0, \forall i\in\mathcal{C}/\mathcal{C}_r \Longrightarrow N_A^i \mu \geq 0, \forall i\in\mathcal{C}_r$). Then, if \eqref{eq:polar_cone_B} is feasible, there exists a minimal representation of $B_A$ in terms of the nullspace bases in $N_A$, i.e., $\nu^\star_i = 0, \forall i\in\mathcal{C}_r$.
\end{proposition}
\begin{proof}
Let $\nu^\star$ be an arbitrary optimal solution of \eqref{eq:polar_cone_B}. First if a row $i\in\mathcal{C}_r$ of $N_A$ is \emph{redundant}, then there exist coefficients
$\alpha_j \ge 0$ such that
\begin{equation}\label{eq:redundancy}
N_{A}^i = \sum_{j\neq i} \alpha_j N_{A}^j.
\end{equation}
Assume, for contradiction, that
\begin{equation}\label{eq:positive_weight}
\nu^\star_i > 0.
\end{equation}
Define a vector $d \in \mathbb{R}^c$ by
\begin{equation}\label{eq:d_def}
d_i := -1,
\qquad
d_j := \alpha_j \ \text{ for } j\neq i.
\end{equation}
Using \eqref{eq:redundancy}, we compute
\[
N_A^\top d
= -N_{A}^i + \sum_{j\neq i} \alpha_j N_{A}^j
= 0.
\]
Hence, for all $t\in\mathbb{R}$,
\begin{equation}\label{eq:feasible_ray}
N_A^\top(\nu^\star + t d) = B_A.
\end{equation}
For any $t \in (0,\nu^\star_i]$, we have
\[
(\nu^\star + t d)_i = \nu^\star_i - t \ge 0,
\qquad
(\nu^\star + t d)_j = \nu^\star_j + t\alpha_j \ge 0,
\]
so $\nu^\star + t d$ is feasible for \eqref{eq:polar_cone_B}.
The change in the objective value is
\begin{equation}\label{eq:objective_change}
1_c^\top(\nu^\star + t d)
= 1_c^\top \nu^\star
+ t\left(\sum_{j\neq i} \alpha_j - 1\right).
\end{equation}
We now distinguish three exhaustive cases:
\noindent\emph{Case 1: $\sum_{j\neq i}\alpha_j < 1$.}
Then \eqref{eq:objective_change} implies
$1_c^\top(\nu^\star + t d) < 1_c^\top \nu^\star$
for any $t>0$, contradicting the optimality of $\nu^\star$.
\noindent\emph{Case 2: $\sum_{j\neq i}\alpha_j > 1$.}
Then $1_c^\top d > 0$. Since $N_A^\top(-d)=0$, for sufficiently small $t>0$
the vector $\nu^\star - t d$ remains nonnegative and feasible, and satisfies
\[
1_c^\top(\nu^\star - t d)
< 1_c^\top \nu^\star,
\]
again contradicting optimality.
\noindent\emph{Case 3: $\sum_{j\neq i}\alpha_j = 1$.}
Then \eqref{eq:objective_change} yields
\[
1_c^\top(\nu^\star + t d) = 1_c^\top \nu^\star
\quad \forall t\in(0,\nu^\star_i].
\]
Choosing $t=\nu^\star_i$, we obtain a feasible solution with the same objective
value and with the $i$-th component equal to zero.
In all cases, the assumption \eqref{eq:positive_weight} leads to a contradiction
with optimality or shows that an alternative optimal solution exists with
$\nu_i=0$. Therefore, there exists an optimal solution of \eqref{eq:polar_cone_B} satisfying
$\nu^\star_i=0$.
\end{proof}
}

{
\color{red}
The following proposition formalizes Remark \ref{rem:B_cone_components}:
\begin{corollary}\label{cor:cone_boundary}
        Assume that \eqref{eq:polar_cone_B} is feasible with minimal solution $\nu^\star \in \mathbb{R}^c$ per Prop. \ref{prop:minimal_sol}. Further consider the cone $\mathcal{N} \triangleq \{ \mu \in {\color{blue}\mathbb{R}^k} | N_A \mu \geq 0\}$, where some of the rows of $N_A$ may be redundant. Let $N_A^i \in \mathbb{R}^{1\times k}, i\in \mathcal{C}$ denote the rows of $N_A$ and also let $\mathcal{C}_r \subset \mathcal{C}$ denote the indices of the redundant rows of $N_A$ (redundancy means that: $N_A^i \mu \geq 0, \forall i\in\mathcal{C}/\mathcal{C}_r \Longrightarrow N_A^i \mu \geq 0, \forall i\in\mathcal{C}_r$). Then, if $\exists i \in \mathcal{C}/\mathcal{C}_r$ such that $\nu^\star_i = 0 $, this implies that $B_A \in \partial \mathcal{N}_c$.   
\end{corollary}
\begin{proof}
    The polar cone $\mathcal{N}_C$ \eqref{eq:polar_cone} can be equivalently expressed by its generator \cite{linopt_book}:
    \begin{equation}
        \mathcal{N}_C
        = 
        \left\{
            -N_A^\top \nu 
            :
            \nu \in \mathbb{R}^c_{\geq 0}  
        \right\}.
    \end{equation}
Accounting for redundancy, the corresponding minimal representation is:
\begin{equation}
     \mathcal{N}_C
        = 
        \left\{
            -\sum_{i\in\mathcal{C}/\mathcal{C}_r}\nu_i(N_A^i)^\top  
            :
            \nu_i \geq 0, \forall i \in \mathcal{C}/\mathcal{C}_r  
        \right\}.
\end{equation}
The boundary of $\mathcal{N}_C$ then is given by (see \cite{rockafellar-1970a}):
\begin{equation}\label{eq:prop:bound:1}
    \begin{gathered}
        \partial\mathcal{N}_C
        = \\\
        \left\{
            -\sum_{i\in\mathcal{C}/\mathcal{C}_r}\nu_i(N_A^i)^\top  
            :
            \exists j\in \mathcal{C}/\mathcal{C}_r   \textrm{ with } \nu_j = 0
            \textrm{ and }
        \right. 
        \\
        \left.
            \nu_i \geq 0, \forall i \in \mathcal{C}/\mathcal{C}_r 
        \right\}.
    \end{gathered}
\end{equation}
However, note that by \eqref{eq:polar_cone_B}:
\begin{equation}
    -N_A^\top \nu^\star = -\sum_{i\in\mathcal{C}}\nu_i^\star (N_A^i)^\top = -B_A.
\end{equation}
However, since $\nu^\star $ is the minimal representation, i.e. $\sum_{i\in\mathcal{C}}\nu_i^\star (N_A^i)^\top = \sum_{i\in\mathcal{C}/\mathcal{C}_r}\nu_i^\star (N_A^i)^\top$:
\begin{equation}\label{eq:prop:bound:2}
     -B_A = -\sum_{i\in\mathcal{C}/\mathcal{C}_r}\nu_i^\star (N_A^i)^\top.
\end{equation}
Therefore, if $\exists i \in \mathcal{C}/\mathcal{C}_r$ such that $\nu^\star_i = 0 $ along with \eqref{eq:prop:bound:2}, $B_A$ satisfies \eqref{eq:prop:bound:1}, and therefore the former lies on the boundary of the polar cone, concluding the proof. 
\end{proof}
}
{
\color{blue}
\begin{remark}
    Solving the LP \eqref{eq:polar_cone_B} does not necessarily guarantee that $\nu^\star$ is the minimal solution. Nevertheless, Cor. \ref{cor:cone_boundary} can be employed to inform which constraints (corresponding to rows of the matrix $N_A$) render $\mathcal{U}$ empty. Since $B_A \in \mathcal{N}_C$ is a necessary and sufficient condition for non-emptiness of $\mathcal{U}$, the smaller-valued elements of $\nu^\star$ \textbf{may} indicate that small variations of the corresponding constraints may cause $B_A$ to exit $\mathcal{N}_C$ thus rendering $\mathcal{U}$ empty. This will be employed in Subsec. \ref{sec:meth:subsec:local_search} for local constraint selection. 
\end{remark}
}

\subsection{Feasibility Analysis under Disregarded Constraints}
In this section we reconsider the case where only a subset of the constraints are enforced, encoded through configurations $P\in\mathcal{P}$ and their associated index sets $\mathcal{C}_P$.
Per Remark \ref{rem:B_cone_components}, the magnitude of the polar cone components of $B_A$, encoded by $\nu^\star$ in \eqref{eq:polar_cone_B} quantify how ``far'' (see the norm defined in \ref{rem:B_cone_components}) $B_A$ is from the boundary of the polar cone. Consider now a configuration $P\in\mathcal{P}_f$. Since some of the constraints have been disregarded in $P$, the following modified LP can be used to assess how far away it is from infeasibility:
\begin{equation}\label{eq:polar_cone_B:star}
    \begin{gathered}
    \nu^\star_{P} \in \underset{\nu\in\mathbb{R}{\color{blue}^c}}{\arg\min}\left\{\left( P - \frac{1}{2}{1}_c \right)^\top \nu \right\}, \\
    \textrm{s.t.: } N_{A}^\top \nu = B_{A}, \ \textrm{diag}\left( P\right) \nu \geq 0,
    \end{gathered}   
    \tag{\ref{eq:polar_cone_B}*}
\end{equation}
Note that for $P = {1}_c$ the equality constraint in \eqref{eq:polar_cone_B:star} are the same as in \eqref{eq:polar_cone_B}. Therefore, the former is feasible iff the latter is feasible as well. Hence, according to Thm. \ref{thm:feas_lp_farkas} and Props. \ref{prop:cone}, \ref{prop:lp_cone_feas}, $\mathcal{U}$ is nonempty iff \eqref{eq:polar_cone_B:star} is feasible. We generalize this in the following theorem:
\begin{theorem}\label{thm:conf_eval}
    A configuration $P\in\mathcal{P}$ is feasible iff \eqref{eq:polar_cone_B:star} is feasible. 
\end{theorem}
\begin{proof}
The LP \eqref{eq:polar_cone_B:star} can be recast as:
    \begin{equation}\label{eq:polar_cone_B:star:2}
    \begin{gathered}
    \nu^\star_{P}  \in \underset{\nu_e\in\mathbb{R}^{n_e},\nu_d\in\mathbb{R}^{n_d}}{\arg\min}\left\{ \frac{1}{2}{1}_{n_e}^\top \nu_e - \frac{1}{2}{1}_{n_d}^\top \nu_d \right\}, \\
    \textrm{s.t.: } N_{A,e}^\top \nu_e +  N_{A,d}^\top \nu_d 
    = B_{A}, \  \nu_e \geq 0,
    \end{gathered}   
    \tag{\ref{eq:polar_cone_B}**}
\end{equation}
where $n_e,n_d \in \mathcal{C}$ such that $n_e + n_d = c$ denote the number of enforced and disregarded constraints, respectively, and $N_{A,e}\in\mathbb{R}^{ n_e \times k}, N_{A,d}\in\mathbb{R}^{ n_d \times k}$ denote the column vectors of the nullspace basis matrix that correspond to the enforced and the disregarded constraints, respectively.
\par 
(Sufficiency:) If \eqref{eq:polar_cone_B:star} (and thus \eqref{eq:polar_cone_B:star:2}) is feasible, then there exists a $\mathbb{R}^{n_e}\ni\nu_e \geq 0$ and $\nu_d \in \mathbb{R}^{n_d}$ such that $N_{A,e}^\top \nu_e +  N_{A,d}^\top \nu_d 
    = B_{A}$. Additionally, let $\nu_d^+ \in \mathbb{R}^{n_d^+},\nu_d^-\in \mathbb{R}^{n_d^-}$ denote vectors with the positive/negative elements of $\nu_d$, with $n_d^+,n_d^- \in\mathbb{N}$ denoting the number of positive/negative elements of $\nu_d$ respectively. Then, w.l.o.g., (by rearranging the rows of $N_A$), the equality constraints in \eqref{eq:polar_cone_B:star:2} can be expressed as:
\begin{equation}\label{eq:polar_cone_B:star:3}
\begin{gathered}
   N_A^\top \nu = B_{A} \Leftrightarrow \\  
    \begin{bmatrix}
         N_{A,e}^\top, & 
        N_{A,d,+}^\top,  &
         N_{A,d,-}^\top 
     \end{bmatrix}
     \begin{bmatrix}
        \nu_e \\ 
        \nu_d^+ \\
        \nu_d^-
     \end{bmatrix} = \\
     \begin{bmatrix}
             N_{A,e}^\top, & 
             N_{A,d,+}^\top,  &
             N_{A,d,-}^\top 
         \end{bmatrix}
          \begin{bmatrix}
             B_{e}^\top, & 
             B_{d,+}^\top,  &
             B_{d,-}^\top 
         \end{bmatrix}^\top,
\end{gathered}
\end{equation}
with $\nu_e \geq 0, \nu_d^+ \geq 0, \nu_d^- < 0$, or equivalently:
\begin{equation}\label{eq:polar_cone_B:star:4}
    \begin{gathered}
        \begin{bmatrix}
             N_{A,e}^\top, & 
             N_{A,d,+}^\top,  &
             -N_{A,d,-}^\top 
         \end{bmatrix}
         \begin{bmatrix}
            \nu_e \\ 
            \nu_d^+ \\
            -\nu_d^-
         \end{bmatrix} = \\
         \begin{bmatrix}
             N_{A,e}^\top, & 
             N_{A,d,+}^\top,  &
             -N_{A,d,-}^\top 
         \end{bmatrix}
          \begin{bmatrix}
             B_{e}^\top, & 
             B_{d,+}^\top,  &
             -B_{d,-}^\top 
         \end{bmatrix}^\top 
          \Leftrightarrow \\
         \overline{N}_A^\top \bar{\nu} = \overline{B}_A,
    \end{gathered}
\end{equation}
with $\nu_e \geq 0, \nu_d^+ \geq 0, -\nu_d^- > 0$ and where $N_{A,e} \in \mathbb{R}^{n_e \times k}, N_{A,d,+} \in \mathbb{R}^{n_d^+ \times k}, N_{A,d,-}  \in \mathbb{R}^{n_d^- \times k}$ are formed by the columns of ${N}_A$ that correspond to $\nu_e, \nu_d^+, \nu_d^- $ while  $B_{e}  \in \mathbb{R}^{n_e}, B_{d,+}\in \mathbb{R}^{n_d^+}, B_{d,-}\in \mathbb{R}^{n_d^-}$ are formed by the elements of $B$ corresponding to the vectors $\nu_e, \nu_d^+, \nu_d^- $. This however implies that $-\overline{B}_A$ lies within the polar cone formed by the matrix $\overline{N}_A$. Note that by negating the matrix $N_{A,d,-}^\top $ and vectors $B_{d,-}, \nu_d^-$ in \eqref{eq:polar_cone_B:star:4}, this is equivalent to imposing the complementary constraints, i.e., the set:
\begin{equation}\label{eq:set_minus}
    \left\{
        u\in\mathbb{R}^m 
        \left|
            \begin{bmatrix}
                A_{e}^\top \\
                A_{d,+}^\top \\
                -A_{d,-}^\top
            \end{bmatrix} u 
            \leq 
            \begin{bmatrix}
                B_{e} \\
                B_{d,+} \\
                -B_{d,-}
            \end{bmatrix}
        \right.
    \right\}.
\end{equation}
{
\color{red}
To see this, note that flipping the sign of any \textbf{row} of $N_A$, is equivalent to flipping the sign of the associated Lagrange Multiplier (see Fig. \ref{fig:theory}) which is exactly equivalent to enforcing the complementary constraint, i.e., note that $\lambda = N_A \mu $ and thus $\lambda \lessgtr 0 \Leftrightarrow  N_A \mu \lessgtr 0 $.
}
Therefore, through Prop. \ref{prop:cone}, the set \eqref{eq:set_minus} is non-empty. Going back to our original definition, the aforementioned set can  be expressed as \eqref{eq:input_set:config}, which renders configuration $P$ feasible. 
\\
(Necessity:) If $P$ is feasible, then through Prop. \ref{prop:cone}, $-B_A$ belongs to the polar cone of $N_A$, which implies that a solution to:
\begin{equation}
         N_{A,e}^\top
        \nu_e  = B_{A}, \nu_e\geq 0,
\end{equation}
exists, and therefore a solution to \eqref{eq:polar_cone_B:star:2} and thus to \eqref{eq:polar_cone_B:star} also exists, concluding the proof. 
\end{proof}



\section{Methods for Configuration Selection}\label{sec:conf_selec_methods}
In this section we provide configuration selection methods based on the preceding analysis.
{\color{blue}
\subsection{Iterative Method for Constraint Selection}
}
First the following feasibility check algorithm is proposed in Alg. \ref{alg:feas_eval}.
\begin{algorithm}
\caption{Feasibility Check: {\fontfamily{qcr}\selectfont FC}$(A,B,P)$}\label{alg:feas_eval}
\begin{algorithmic}[1]
\State  Given $A\in{\color{blue}\mathbb{R}^{m\times c}},B\in{\color{blue}\mathbb{R}^{c}}$, configuration $P\in\mathcal{P}$,
\State  Compute nullspace basis $N_A\in{\color{blue}\mathbb{R}^{c\times k}}$, $B_A= N_A^\top B$,
\State  Solve LP:
\begin{equation}\label{eq:fc:alg}
    \begin{gathered}
    \nu \in \underset{\nu\in\mathbb{R}{\color{blue}^c}}{\arg\min}\left\{\left( P - \frac{1}{2}{1}_c \right)^\top \nu \right\}, \\
    \textrm{s.t.: } N_{A}^\top \nu = B_{A}, \ \textrm{diag}\left( P\right) \nu \geq 0,
    \end{gathered}   
\end{equation}
\If{Problem \eqref{eq:fc:alg} is feasible}
    \State \Return TRUE, $\nu$
\ElsIf{Problem \eqref{eq:fc:alg} is infeasible}
    \State \Return FALSE
\EndIf
\end{algorithmic}
\end{algorithm}
\par 
An important observation on the LP \eqref{eq:polar_cone_B:star:2} in the context of the proof of Thm. \ref{thm:conf_eval} is that owing to minimizing the magnitude of (the positive elements of) $\nu_e\in\mathbb{R}^{n_e}_{\geq 0}$, some of the elements of the disregarded constraints' vector $\nu_d$ may become positive (i.e., $\nu_d^+ $). The is an effect of minimizing the negated sum of the elements through the term $- \frac{1}{2}{1}_{n_d}^\top \nu_d$ in \eqref{eq:polar_cone_B:star:2}; decreasing the magnitude of the elements of $\nu_e$ (accomplished by minimizing $\frac{1}{2}{1}_{n_e}^\top \nu_e$) allows for some of the elements of $\nu_d$ to become positive. This is formalized in the following corollary:
\begin{corollary}\label{col:adding_cons}
Given a feasible configuration $P\in\mathcal{P}_f$, consider the LP \eqref{eq:polar_cone_B:star} and its equivalent \eqref{eq:polar_cone_B:star:2}. Denote its minimizer as $\nu^\star = \left[ \nu_e ^\top, (\nu_d^+)^\top, (\nu_d^-)^\top \right]^\top$, where $\nu_e$ corresponds to the enforced constraints and $\nu_d^+ \in \mathbb{R}^{n_d^+}_{\geq 0}, \nu_d^- \in \mathbb{R}^{n_d^-}_{< 0}$ correspond to the disregarded constraints, with the associated constraint matrices/vectors:
\begin{equation}
   \mathcal{U} =  \left\{
        u\in\mathbb{R}^m 
        \left|
            \begin{bmatrix}
                A_{e}^\top \\
                A_{d,+}^\top \\
                A_{d,-}^\top
            \end{bmatrix} u 
            \leq 
            \begin{bmatrix}
                B_{e} \\
                B_{d,+} \\
                B_{d,-}
            \end{bmatrix}
        \right.
    \right\},
\end{equation}
where $A_{e} \in \mathbb{R}^{m \times n_e}, A_{d,+} \in \mathbb{R}^{m \times  n_d^+ }, A_{d,-}  \in \mathbb{R}^{m \times n_d^-}$ are formed by the columns of $A$ that correspond to $\nu_e, \nu_d^+, \nu_d^- $ while  $B_{e}  \in \mathbb{R}^{n_e}, B_{d,+}\in \mathbb{R}^{n_d^+}, B_{d,-}\in \mathbb{R}^{n_d^-}$ are formed by the elements of $B$ corresponding to the vectors $\nu_e, \nu_d^+, \nu_d^- $. Then, enforcing the constraints 
\begin{equation}
   \mathcal{U} =  \left\{
        u\in\mathbb{R}^m 
        \left|
            \begin{bmatrix}
                A_{e}^\top \\
                A_{d,+}^\top 
            \end{bmatrix} u 
            \leq 
            \begin{bmatrix}
                B_{e} \\
                B_{d,+} 
            \end{bmatrix}
        \right.
    \right\},
\end{equation}
yiels a feasible problem. 
% Enforcing the constraints corresponding to $\nu_d^+$ yields a feasible set of constraints. 
\end{corollary}
\begin{proof}
    Following the arguments of the proof of Thm. \ref{thm:conf_eval}, Eq. \eqref{eq:polar_cone_B:star:4}, positivity of $\nu_e, \nu_d^+$ implies that the vector $-\left[  B_{e}, B_{d,+}  \right]^\top$ lies within the polar cone associated with the matrix $ \begin{bmatrix}
             N_{A,e}^\top, & 
             N_{A,d,+}^\top 
         \end{bmatrix}$. Therefore, through Prop. \ref{prop:cone} enforcing the aforementioned constraints yields a feasible set.   
\end{proof}
Cor. \ref{col:adding_cons} can hence be employed to formulate an iterative method for enforcing more constraints. In Alg. \ref{alg:iter}, starting from a feasible configuration, the LP \eqref{eq:feas_lp:star} is solved and the configuration is updated by enforcing the disregarded constraints that correspond to positive components of the LP solution in \eqref{eq:feas_lp:star}. Owing to Cor. \ref{col:adding_cons}, each of the newly enforced constraints will yield feasible constraint sets.  
\par 
However, Alg. \ref{alg:iter} does not necessarily yield the maximum number of feasible constraints; note that an initial feasible configuration defines a polytope within $\mathbb{R}^m$. Since we have demonstrated that Alg. \ref{alg:iter} adds feasible constraints, its iterations can  only produce subsets within the initial polytope. Any other sets disjoint to the initial configuration's polytope, {\color{red}if they exist,} cannot be obtained through this algorithm. 
\begin{algorithm}
\caption{Iterative Constraint Addition: {\fontfamily{qcr}\selectfont ICA}$(A,B,P^{(0)})$}\label{alg:iter}
\begin{algorithmic}[1]
\State  Given $A\in{\color{blue}\mathbb{R}^{m\times c}},B\in{\color{blue}\mathbb{R}^{c}}$, initial configuration $P^{(0)}\in\mathcal{P}$,
\State Check feasibility: $\mathrm{FLAG},\nu \gets${\fontfamily{qcr}\selectfont FC}$(A,B,P^{(0)})$,
\If{FLAG == FALSE}
    \State Set $P^{(0)} \gets \overline{P}$ according to the hard constraints:
    \begin{equation}
    \overline{P} = \begin{cases}
         1, & i \in \mathcal{C}_h\\
         0, & i \in \mathcal{C}_s
    \end{cases}
    \notag 
    \end{equation}
    \State Get $\nu^{(0)}$: $\mathrm{FLAG},\nu^{(0)} \gets${\fontfamily{qcr}\selectfont FC}$(A,B,P^{(0)})$,
\EndIf
\State  Split $\nu^{(0)}$ into $\left[ \nu_{e,0}^\top, (\nu_{d,0}^{+})^\top, (\nu_{d,0}^{-})^\top\right]^\top$, set $i\gets 0$,
\While{$\nu_{d,i}^{+}$ is not empty}
    \State  Update configuration:
    \begin{equation}
        P^{(i+1)}_j = 
        \begin{cases}
            1,& \nu^{(i)}_j \geq 0 \\
            0,& \nu^{(i)}_j < 0
        \end{cases}, \forall j \in \{1,\cdots,C\},
    \end{equation}
    \State  Solve LP:
    \begin{equation}
        \begin{gathered}
        \nu^{(i+1)} \in \underset{\nu\in\mathbb{R}{\color{blue}^c}}{\arg\min}\left\{\left( P^{(i+1)} - \frac{1}{2}{1}_c \right)^\top \nu \right\}, \\
        \textrm{s.t.: } N_{A}^\top \nu = B_{A}, \ \textrm{diag}\left( P^{(i+1)}\right) \nu \geq 0,
        \end{gathered}   
    \end{equation}
    \State  Split $\nu^{(i+1)}$ into $\left[ \nu_{e,i+1}^\top, (\nu_{d,i+1}^{+})^\top, (\nu_{d,i+1}^{-})^\top\right]^\top$,
    \State  $i\gets i+1$,
\EndWhile
\State \Return $P^{(i)},\nu^{(i)}$. 
\end{algorithmic}
\end{algorithm}
\par 
Alg. \ref{alg:iter} requires initialization with a feasible configuration. Note the equality constraint of \eqref{eq:polar_cone_B:star} yields:
\begin{equation}
    \begin{gathered}
        N_A^\top \nu = B_A = N_A^\top B \Leftrightarrow
        N_A^\top \left( \nu - B \right) = 0. 
    \end{gathered}
\end{equation}
Hence, for $\nu = B$ the linear constraints of the feasibility LP \eqref{eq:polar_cone_B:star} are satisfied and it is feasible as result, rendering the corresponding configuration feasible. Enforcing the constraints for which $\nu = B \geq 0$, can serve as an initialization for Alg. \ref{alg:iter}. 

\subsection{Local Search \& Ranking Constraints for Constraint Selection}\label{sec:meth:subsec:local_search}
Consider again the LP \eqref{eq:polar_cone_B:star}. \textcolor{orange}{As noted in Rem.~\ref{rem:B_cone_components}, the magnitude of each component of $\nu_P^*\in\mathbb{R}^c$ in~\eqref{eq:polar_cone_B:star} quantifies} the impact of each constraint towards feasibility/infeasibility, since it corresponds to a component of the vector $B_A$ within the polar cone of the matrix $N_A$. We remind the reader that $B_A$ lying within the aforementioned polar cone is a necessary and sufficient condition for the corresponding set of constraints to be feasible (see Prop. \ref{prop:cone}). Therefore, given a configuration $P\in\mathcal{P}$, the constraints can be sorted based on their corresponding value of the solution to the LP \eqref{eq:polar_cone_B:star} denoted by $\nu^\star_P$. For enforced constraints, the smaller their corresponding (positive) element of $\nu^\star_P\in \mathbb R^c$, the ``closer'' the former is to rendering the problem infeasible (see Cor. \ref{cor:cone_boundary}, \ref{cor:cone_boundary:2}). For disregarded constraints, the larger the corresponding (negative) element of $\nu^\star_P$ is (i.e., the closer it is to zero), the more likely it is that the constraint is feasible (see Cor. \ref{cor:cone_boundary}, \ref{cor:cone_boundary:2}). This is formalized in the following Corollary:
{
\color{red}
\begin{corollary}\label{cor:cone_boundary:2}
        Assume that \eqref{eq:polar_cone_B:star} is feasible with minimal solution $\nu^\star \in \mathbb{R}^c$ per Prop. \ref{prop:minimal_sol}. Further consider the cone $\mathcal{N} \triangleq \{ \mu \in {\color{blue}\mathbb{R}^k} | N_A \mu \geq 0\}$, where some of the rows of $N_A$ may be redundant. Let $N_A^i \in \mathbb{R}^{1\times k}, i\in \mathcal{C}$ denote the rows of $N_A$ and also let $\mathcal{C}_r \subset \mathcal{C}$ denote the indices of the redundant rows of $N_A$ (redundancy means that: $N_A^i \mu \geq 0, \forall i\in\mathcal{C}/\mathcal{C}_r \Longrightarrow N_A^i \mu \geq 0, \forall i\in\mathcal{C}_r$). Then, if $\exists i \in \mathcal{C}/\mathcal{C}_r$ such that $\nu^\star_i = 0 $, this implies that $B_A \in \partial \mathcal{N}_c$.   
\end{corollary}
\begin{proof}
    The proof follows along the same lines as Prop. \ref{prop:minimal_sol}, Cor. \ref{cor:cone_boundary}, through a change of variables similar to Thm. \ref{thm:conf_eval}. More specifically, recasting  \eqref{eq:polar_cone_B:star} as \eqref{eq:polar_cone_B:star:2} and setting $\nu_d = -\nu_d'$ yields the original LP \eqref{eq:polar_cone_B}, and Prop. \ref{prop:minimal_sol}, Cor. \ref{cor:cone_boundary} then apply. 
\end{proof}
}
Finally, based on Cor. \ref{cor:cone_boundary:2} Alg. \ref{cor:cone_boundary:2} is proposed for local constraint selection. 
% \par 
{\color{blue}
% \begin{proposition}
    
% \end{proposition}

% The above are implied from Thm. \ref{thm:conf_eval}. Existence of a solution $\nu^\star$ corresponds to expressing the vector $B_A = N_A^\top B $ as a linear combination of the null-space basis vectors, which in turn corresponds to $B_A$ lying within the polar cone formed by the matrix $N_A$. 
% \par 
% To see this, consider the matrices $N_A\in\mathbb{R}^{c \times k}, B_A \in \mathbb{R}^k$ and configurations $P_1,P_2\in\mathcal{P}$ and their variations $\Delta_N \in\mathbb{R}^{c \times k}, \Delta_B \in \mathbb{R}^k, N_A' \triangleq N_A + \Delta_N $ and $B_A'\triangleq B_A + \Delta_B$. Assume w.l.o.g. that:
% \begin{equation}\label{eq:constr:rank:1}
%     \exists \mathbb{R}^c \ni\nu_1 :\textrm{diag}(P_1)\nu_1\geq 0, N_A^\top \nu_1 = B_A,
% \end{equation}
% such that $\mathbb{R}^c \ni \nu_1 = \left[ (\nu_1^{\pm})^\top,  (\nu_1^0)^\top \right]^\top$, where the elements of $\nu_1^{\pm}$ are strictly positive/negative and the elements of $\nu_1^{0}$ are equal to zero, which implies that $P_1 \in \mathcal{P}_f$ for the matrices $N_A,B_A$. Further assume that $P_1$ is infeasible w.r.t. the matrices $N_A',B_A'$. We will show that if $P_2$ is feasible for the matrices $N_A',B_A'$, i.e.,:
% \begin{equation}\label{eq:constr:rank:2}
%     \exists \mathbb{R}^c \ni\nu_2 :\textrm{diag}(P_2)\nu_2\geq 0, (N_A')^\top \nu_2 = B_A',
% \end{equation}
% this then necessarily implies that $\nu_2$ has the following form: $\mathbb{R}^c \ni \nu_2=  \left[ (\nu_1^{\pm})^\top,  (\nu_2^-)^\top \right]^\top$ with the elements of $\nu_2^-$ being strictly negative.
% \par 
% To see this, we begin by assuming that \eqref{eq:constr:rank:1}, \eqref{eq:constr:rank:2} hold, and express $\nu_2 = \left[ \nu_{2,1}^\top, \nu_{2,2}^\top\right]^\top$, where $\textrm{dim}(\nu_{2,1}) = \textrm{dim}(\nu_{1}^{\pm})$ and $\textrm{dim}(\nu_{2,2}) = \textrm{dim}(\nu_{1}^{0})$.
% From \eqref{eq:constr:rank:2}:
% \begin{equation}\label{eq:constr:rank:3}
%     \begin{gathered}
%         (N_A')^\top \nu_2 = B_A' \Leftrightarrow \\
%         (N_A + \Delta_N)^\top \nu_2= B_A + \Delta_B \overset{\eqref{eq:constr:rank:1}}{\Leftrightarrow} \\
%         (N_A + \Delta_N)^\top \nu_2= N_A^\top \nu_1 + \Delta_B \Leftrightarrow \\
%         N_A^\top \left( \nu_2 - \nu_1\right)
%         + \Delta_N^\top \nu_2 = \Delta_B.
%     \end{gathered}
% \end{equation}
% \begin{figure}
%     \centering
%     \includegraphics[trim={2cm 2.5cm 10cm 4cm},clip,width=0.5\textwidth]{figures/theory_2.pdf}
%     \caption{Visualization of the variations on $N_A,B_A$.}
%     \label{fig:constr:ranking:vis}
% \end{figure}
% Note now that for the variations $\Delta_N,\Delta_B$ to render $P_1$ (which is feasible w.r.t. the matrices $N_A,B_A$) infeasible for the matrices $N_A',B_A'$, these variations should alter the constraints corresponding to $\nu_1^0$, and thus the corresponding columns/elements of $N_A,B_A$. Therefore, splitting $N_A = \left[ (N_A^{\pm})^\top,(N_A^0)^\top \right]^\top$ according to the negative/non-negative elements of $\nu_1$,  this implies for small $\Delta_N, \Delta_B$ (see Fig. \ref{fig:constr:ranking:vis}):
% \begin{equation}\label{eq:constr:rank:4}
% \begin{gathered}
%         \Delta_N^\top = \left[ \mathbf{0}, (\Delta_N^0)^\top \right],\quad 
%         \Delta_B^\top = \left[ \mathbf{0}, (\Delta_B^0)^\top \right]^\top \Rightarrow \\
%         \exists \nu_{2,2}<0: (N_A^0+ \Delta_N^0)^\top \nu_{2,2} \approx \Delta_B^0. 
% \end{gathered}
% \end{equation}
% Therefore from \eqref{eq:constr:rank:2}:
% \begin{equation}
%    \begin{gathered}
%         N_A^\top 
%         \begin{bmatrix}
%             \nu_{2,1}- \nu_1^{\pm}
%             \\
%             \nu_{2,2}
%         \end{bmatrix}
%         + \Delta_N^\top
%         \begin{bmatrix}
%             \nu_{2,1}
%             \\
%             \nu_{2,2}
%         \end{bmatrix}
%         = \Delta_B {\Leftrightarrow} \\
%         \left[ (N_A^{\pm})^\top,(N_A^0)^\top \right]
%         \begin{bmatrix}
%             \nu_{2,1}- \nu_1^{\pm}
%             \\
%             \nu_{2,2}
%         \end{bmatrix}
%         + 
%         \left[ \mathbf{0}, (\Delta_N^0)^\top \right]
%         \begin{bmatrix}
%             \nu_{2,1}
%             \\
%             \nu_{2,2}
%         \end{bmatrix}
%         = 
%         \begin{bmatrix}
%              \mathbf{0} \\ \Delta_B^0
%         \end{bmatrix}.
%         \notag 
%    \end{gathered}
% \end{equation}
% Selecting $\nu_{2,1} = \nu_1^\pm$, the above yields:
% \begin{equation}
%     \begin{gathered}
%                 (N_A^0 + \Delta_N^0)^\top\nu_{2,2} =  \Delta_B^0,
%     \end{gathered}
% \end{equation}
% which is approximately satisfied for $\nu_{2,2}<0$ per \eqref{eq:constr:rank:4}, and therefore $\nu_2=  \left[ \nu_{2,1}^\top,  \nu_{2,2}^\top \right]^\top=  \left[ (\nu_1^{\pm})^\top,  (\nu_2^-)^\top \right]^\top$. 
% \par 
% The above imply that infeasibility under small variations, can be addressed by disregarding only the constraints corresponding to the zero-valued elements in the solution of \eqref{eq:polar_cone_B:star}. Hence, enforced constraints corresponding to the smallest values in the solution to \eqref{eq:polar_cone_B:star} are more likely to cause the corresponding problem to become infeasible. 
% \par
% The above are not in the form of a theorem/proof because we aim on exploring such variational approaches more formally in the future. 
% In the context of constraint selection, this constraint ranking can be used to perform a local configuration search, presented in Alg. \ref{alg:local_search}.
% \begin{remark}
%     The variations considered above may encode changes in the constraint matrices $A(x,t), B(x,t)$ over time, or due to the evolution of System \eqref{eq:dynamics}'s state $x(t)$ (both of which alter the nullspace bases' matrix $N_A$ and vector $B_A = N_A^\top B$). Analyzing the properties of such variations and their effect on constraint selection is left as part of future research. 
% \end{remark}
}




\begin{algorithm}
\caption{Local Configuration Search: {\fontfamily{qcr}\selectfont LCS}$(A,B,P,P^{-},\nu^{-},D)$}\label{alg:local_search}
\begin{algorithmic}[1]
\State  Given $A\in{\color{blue}\mathbb{R}^{m\times c}},B\in{\color{blue}\mathbb{R}^{c}}$, current configuration $P\in\mathcal{P}$, last feasible configuration $P^{-}\in\mathcal{P}_f$, LP solution vector \eqref{eq:polar_cone_B:star} $\nu^-$, search depth $\mathbb{N}\ni D\leq c$,
\State Check feasibility: $\mathrm{FLAG},\nu \gets${\fontfamily{qcr}\selectfont FC}$(A,B,P)$,
\State Initialize hard constraints configuration:
    \begin{equation}
    \overline{P}_i = \begin{cases}
         1, & i \in \mathcal{C}_h\\
         0, & i \in \mathcal{C}_s
    \end{cases},
    \notag 
\end{equation}
\If{ FLAG  = TRUE}
    \State Enforce additional constraints: 
     \begin{equation}
         P^+,\nu^+ \gets \textrm{\fontfamily{qcr}\selectfont ICA}(A,B,P),\ \notag 
     \end{equation}
     \State Save temporary variable: $P^{t}\gets P^+$,
     \For{$i = 1:D$}
        \State Split $\nu^+$ into $\left[ \nu_{e}^\top, \nu_{d}^\top\right]^\top$,
        \State 
        \parbox[t]{\dimexpr\linewidth-\algorithmicindent-20pt}{Rank the elements of $\nu_{d}$ into the vector $\nu_{d}^r \in \mathbb{R}^{c_d}_{\leq 0}$ and save the indices of the elements of $\nu_{d}^r$ in $\nu_{d}$ as $\{ i_1,i_2,\cdots ,i_{c_d}\}$},
        \For{$j = 1:c_d$}
            \State Update $P^{t}_{i_j}\gets 1$,
            \State Check feasibility:
            \begin{equation}
                \mathrm{FLAG},\nu^t \gets\textrm{\fontfamily{qcr}\selectfont FC}(A,B,P^t),\notag
            \end{equation}
            \If{FLAG = FALSE}
                \State $P^{t}_{i_j}\gets 0$
            \EndIf
        \EndFor
        \State Update $P^{+} \gets P^t,\ \nu^+ \gets \nu^t$,
     \EndFor
\ElsIf{FLAG = FALSE}
\State Save temporary variable: $P^{t}\gets P$,
    \For{$i = 1:D$}
        \State Split $\nu^-$ into $\left[ \nu_{e}^\top, (\nu_{d}^+)^\top, (\nu_{d}^-)^\top\right]^\top$,
        \State \parbox[t]{\dimexpr\linewidth-\algorithmicindent-20pt}{Rank (the first element is the minimum) the elements of $\nu_{e}$ into the vector $\nu_{e} \in \mathbb{R}^{c_e}_{\geq 0}$ and save the indices of the elements of $\nu_{e}^r$ in $\nu_{e}$ as $\{ i_1,i_2,\cdots ,i_{c_d}\}$},
        \For{$j = 1:c_d$}
            \State Update $P^{t}_{i_j}\gets 0$,
            \State Check feasibility:
            \begin{equation}
                \mathrm{FLAG2},\nu^t \gets\textrm{\fontfamily{qcr}\selectfont FC}(A,B,P^t),\notag
            \end{equation}
            \If{FLAG2 = TRUE}
                \State $\nu^- \gets \nu^t$,
                \State \textbf{BREAK}
            \EndIf
        \EndFor
    \EndFor
    \If{FLAG2 = FALSE}
        \State\parbox[t]{\dimexpr\linewidth-\algorithmicindent-20pt}{
        \begin{equation}
             P^+,\nu^+ \gets \textrm{\fontfamily{qcr}\selectfont ICA}(A,B,\bar{P}),\ \notag 
         \end{equation}
         where $\bar{P}$ is the configuration with only the hard constraints enforced,}
     \EndIf
\EndIf
\State \Return $P^{+}$. 
\end{algorithmic}
\end{algorithm}

\section{Simulations}\label{sec:sims}
In this section, the elements of the previous sections are incorporated into a scheme for efficient feasible configuration search. Herein we consider the case where the constraints are composed of hard and soft constraints indexed by $\mathcal{C}_h,\mathcal{C}_h$ respectively. 


% \subsection{Offline Feasible Configuration Search}


% {\color{red}
% Create this tree-like search structure:
% \begin{itemize}
%     \item Start with tightest shell,
%     \item Activate the feasible rest one by one, ranked by their $\nu$ magnitude and for each of them run Alg. \ref{alg:iter}, save each feasible config. 
%     \item Do shell expansion (remove one constraint) based on minimum magnitude of $\nu$
%     \item The saved feasible configs will still be feasible, do some sort of local search, try to not miss solutions (see Fig. \ref{fig:shells}) this has to do with disregarded constraints that are not part of the saved configs 
% \end{itemize}
% }

\subsection{Online Feasible Constraint Selection}
In this subsection the observation in Rem. \ref{rem:B_cone_components} is employed for online constraint selection. For System \eqref{eq:dynamics} consider the following hybrid system extension:
\begin{equation}\label{eq:dynamics:hybrid}
    \begin{split}
        \mathcal{H}
        \begin{cases}
            \begin{bmatrix}
                \dot{x} \\
                \dot{P}
            \end{bmatrix} = 
            \begin{bmatrix}
                f(x) + g(x)u \\
                {0}_C
            \end{bmatrix}, 
            & 
            x\in\mathbb{R}^n, u \in \mathcal{U}'(P),P\in\mathcal{P}_f
            \\ \\
            \begin{bmatrix}
                {x}^+ \\
                {P}^+
            \end{bmatrix} = 
            \begin{bmatrix}
                x \\
                V
            \end{bmatrix},
            & 
            x\in\mathbb{R}^n, P\notin\mathcal{P}_f, V \in \mathcal{P}_f
        \end{cases}
    \end{split}.
\end{equation}
System \eqref{eq:dynamics:hybrid} models the system evolving continuously under \eqref{eq:dynamics} under a feasible configuration $P\in\mathcal{P}_f$ \eqref{eq:feas_conf_set_def}. When infeasibility occurs, the virtual input $V\in\mathcal{P}_f$ models a change in the applied configuration denoted by $P^+$. Then the system continues to evolve according to the continuous dynamics. Let $\mathcal{T} \triangleq \bigcup_{j=0}^{J} \left([t_j,t_{j+1}] \times\{j\} \right)$ denote the hybrid time, where $j\in\{0,1,\cdots,J\},J\in\mathbb{N}$ denotes the discrete transition index and counts the number of ``jumps'' in the evolution of \eqref{eq:dynamics:hybrid}. For convenience, let any $\mathcal{T}\ni T = (t,j)$ where $t\in[t_j,t_{j+1}],j\in\{0,1,\cdots,J\}$. Then solutions to \eqref{eq:dynamics:hybrid} can be cast as hybrid arcs $x:\mathcal{T}\rightarrow\mathbb{R}^n, P:\mathcal{T}\rightarrow\mathcal{P}, u:\mathcal{T}\rightarrow\mathbb{R}^m$ with an overloading of notation. 
\par 
Consequently, a decision-making scheme for $V\in\mathcal{P}_f$ needs to be designed. When infeasibility occurs, this essentially boils down to starting from the current configuration $P(T)\notin\mathcal{P}_f$ at time $T\in\mathcal{T}$ and state $x(T)$ and finding a feasible one. Formally, we model this search method as a mapping $V:\mathcal{T}\times \mathbb{R}^n \times \mathcal{P} \rightarrow \mathcal{P}_f$. Furthermore, in order to control system \eqref{eq:dynamics} an optimization-based controller can be employed. For instance, consider a QP controller:
    \begin{equation}\begin{gathered}
       \pi_{A,B}(P,x,t)  = \underset{u\in\mathbb{R}^m}{\arg\min}\left\{ \|u - u_{ref}\|_2^2 \right\}\\
        \textrm{s.t.} \quad \textrm{diag}(P)A(x,t)^\top u \leq\textrm{diag}(P)B(x,t),
        \label{eq:cbf-qp}
    \end{gathered}\end{equation}
where $u_{ref}\in\mathbb{R}^m$ denotes a reference input to system~\eqref{eq:dynamics}.
We also include the polar cone components' \eqref{eq:polar_cone_B:star} evolution as the mapping $\nu^\star_P:\mathbb{R}^n\times \mathbb{R}_{\geq 0}\rightarrow\mathbb{R}^c$:
\begin{equation}
    u = \pi_{A,B}(P,x,t),
\end{equation}
\begin{equation}\label{eq:polar_cone_B:star_time}
    \begin{gathered}
    \nu^\star_{P}(x,t)=  \underset{\nu\in\mathbb{R}{\color{blue}^c}}{\arg\min}\left\{\left( P - \frac{1}{2}{1}_c \right)^\top \nu \right\}, \\
    \textrm{s.t.: } N_{A}^\top (x,t)\nu = B_{A}(x,t), \ \textrm{diag}\left( P\right) \nu \geq 0,
    \end{gathered}   
\end{equation}
where $N_A : \mathbb{R}^n\times \mathbb{R}_{\geq 0}\rightarrow{\color{blue}\mathbb{R}^{c\times k}}$ denotes a nullspace basis of the matrix $A(x,t)$ and $B_A(x,t) = N^\top_A(x,t)B(x,t)$, with an overloading of notation. 
We employ two methods to achieve this task.

{\color{red}
\subsubsection{\textbf{Employing Alg. \ref{alg:iter}}}
Alg. \ref{alg:iter} can be used to compute feasible configurations online, which requires an initial feasible configuration as an initialization. Nevertheless, owing to Assum. \ref{assum:hard_const_compatible} the hard constraints' set can be used to this end. 
% This process is described in Alg. \ref{alg:feas_perm_srch:1}, where the latter is called whenever infeasibility is detected. 
% \begin{algorithm}
% \caption{Control with Feasible Configuration Search}\label{alg:feas_perm_srch:1}
% \begin{algorithmic}[1]
% \State
%  \parbox[t]{\dimexpr\linewidth-\algorithmicindent-20pt}{%
% Given the constraint set \eqref{eq:input_set:config} composed of hard and soft constraints indexed by $\mathcal{C}_h,\mathcal{C}_s$, System \eqref{eq:dynamics:hybrid}, an initial hybrid state $\bar{x} \triangleq  x(t_0,0) = x(0,0)\in\mathbb{R}^n$:}
% \State 
%  \parbox[t]{\dimexpr\linewidth-\algorithmicindent-20pt}{%
% Set $P(t_0,0) = P(0,0)  = \overline{P}$ according to the hard constraints:}
% \begin{equation}
%     \overline{P}_i = \begin{cases}
%          1, & i \in \mathcal{C}_h\\
%          0, & i \in \mathcal{C}_s
%     \end{cases}
%     \notag 
% \end{equation}
% \State  Enforce additional constraints through Alg. \ref{alg:iter}:\\
%  \begin{equation}
%      P{(0,0)} \gets \textrm{\fontfamily{qcr}\selectfont ICA} (A(\bar{x},0),B(\bar{x},0),P{(0,0)}),
%      \notag 
%  \end{equation}
%  \State  Set $j\gets 0$,
% \While{System has not converged}
%     \State 
%      \parbox[t]{\dimexpr\linewidth-\algorithmicindent-20pt}{%
%     Evolve System \eqref{eq:dynamics:hybrid} from the state $x(t_j,j)$ and time $t_j$, under the input $\pi_{\left( A,B\right)}(P,x,t),$}
%     \If{ Prob. \eqref{eq:polar_cone_B:star_time} is infeasible for $P(t_j,j)$ }
%         \State  Enforce additional constraints through Alg. \ref{alg:iter}:
%          \begin{equation}
%             \begin{gathered}
%                 P_{j} \gets  
%                  \textrm{\fontfamily{qcr}\selectfont ICA} (A_j,B_j,P_{j}), 
%             \end{gathered}\notag 
%          \end{equation}
%          \State Or Alg. \ref{alg:local_search}
%          \begin{equation}
%             \begin{gathered}
%                 P_{j} \gets 
%                  \textrm{\fontfamily{qcr}\selectfont LCS}
%                  (A_j,B_J,P_{j},P_{j-1},D), 
%             \end{gathered}\notag 
%          \end{equation}
%          \State \parbox[t]{\dimexpr\linewidth-\algorithmicindent-20pt}{
%          where $A_j = A(x(t_j,j),t_j)$,$B_j = B(x(t_j,j),t_j)$,$P_{j} = P{(t_{j},j)}$,$P_{j-1} = P{(t_{j},{j-1})}$,}
%         \State  Set $j\gets j+1$,
%     \EndIf
% \EndWhile
% \end{algorithmic}
% \end{algorithm}


\subsubsection{\textbf{Employing Alg. \ref{alg:local_search}}}
In order to acquire a feasible configuration, one/some of the enforced constraints needs to be disregarded and we remind the reader that enforced constraints correspond to positive values in the corresponding elements of the solution to \eqref{eq:polar_cone_B:star}, denoted by $\nu_P^\star$. Per Rem. \ref{rem:B_cone_components}, upon the solution of \eqref{eq:polar_cone_B:star} the constraints that correspond to the smallest (positive) elements of $\nu_P^\star$ are most likely to cause infeasibility. 
% {\color{red}
% This is implied from Prop. \ref{prop:lp_cone_feas}. Existence of a solution $\nu^\star$ corresponds to expressing the vector $B_A = N_A^\top B $ as a linear combination of the null-space basis vectors, which in turn corresponds to $B_A$ lying within the polar cone formed by the matrix $N_A$. }
All of these elements are combined into Alg. \ref{alg:local_search}.
}

\subsection{Simulation Setup}
 Consider a model for the motion of a robot as:
% \begin{equation}
    \( \dot x(t) = u(t) \),
% \end{equation}
where the control inputs are bounded within the set $U = [-1,1]\times [-1,1] \subset \mathbb R^2$ m/s. The robot starts at the initial position $\bar x = x(t_0)$, where $t_0 \geq \mathbb R_{\geq 0}$, and is tasked to reach a goal position $x_g$ within $T=30$ seconds, while avoiding a set of undesired zones $\mathcal C_s=\{1,\dots, n_s\}$. Each zone $i$ is modeled as a circular disk in $\mathbb R^2$ of radius $r = 1.5$ m, with centers $y_i\in \mathbb R^2$. The zones are divided into static $\mathcal C_N$ and dynamic ones $\mathcal C_D$, where $\mathcal C_N\cap \mathcal C_D=\emptyset$ and $\mathcal C_N\cup \mathcal C_D=\mathcal C_s$. Each dynamic zone $j\in \mathcal C_D$ moves with a known constant velocity $v_j \in \mathbb R^2$, maintaining constant directions and speed. 
\par 
We assume that the robot knows the locations and velocities of the zones. The robot is tasked with reaching the goal within the provided time interval, and with a minimal number of disregarded constraints. Avoiding zone $i \in \mathcal C_s$ is expressed as the superlevel set of the constraint function: 
% \begin{equation\)}
     \(  h_i(x)= ||x-y_i||_2^2 - r^2 \).
% \end{equation}
Furthermore, the distance to the goal is expressed via the Control Lyapunov Function (CLF) constraint:
% \begin{equation}
 \(  V(x)=(x-x_g)^2 \).
% \end{equation}

Using the Control Barrier Function (CBF) and CLF conditions~\cite{ames2016control}, we enforce robot to satisfy the constraints. We then define matrices to compactly represent the constraints. We first define the constraint matrices for the static zones: $A_N(x,t) = \begin{bmatrix}
    A_i(x,t)
\end{bmatrix}_{i \in \mathcal C_N}$ and $B_N(x,t) = \begin{bmatrix}
    B_i(x,t)
\end{bmatrix}_{i \in \mathcal C_N}$, where $A^\top_i(x,t) = -\frac {\partial h_i}{\partial x}$ and $B_i(x,t) = \alpha_i(h_i(x))$ for $i\in \mathcal C_N$. Similarly, we define the constraint matrices for the dynamic zones: $A_D(x,t) = \begin{bmatrix}
    A_j(x,t)
\end{bmatrix}_{j \in \mathcal C_D}$ and $B_D(x,t) = \begin{bmatrix}
    B_j(x,t)
\end{bmatrix}_{j \in \mathcal C_D}$, where $A^\top_j(x,t) = -\frac {\partial h_j}{\partial x}$ and $B_j(x,t) = \alpha_j(h_j(x))+\frac {\partial h_j}{\partial y_j} v_j$ for $j\in \mathcal C_D$. Lastly, we define the matrices for the hard constraints: $A_H(x,t)=\begin{bmatrix}
    \frac {\partial V}{\partial x}^\top & I_m & -I_m
\end{bmatrix}$ and $B_H(x,t) = \begin{bmatrix}
     -\alpha(V(x)) & \mathbf 1_m^\top& \mathbf 1_m^\top
\end{bmatrix}$. Thus, we design a controller in the form of~\eqref{eq:cbf-qp} with
where $A(x,t) = [A_N(x,t) ,A_D(x,t), A_H(x,t) ]$ and $B(x,t) = [B_N(x,t), B_D(x,t), B_H(x,t)]^\top$. The zone avoidance constraints are treated as soft constraints - they may be violated if necessary during navigation - whereas the input bounds and CLF condition are treated as hard constraints that must always be satisfied.


The initial positions of the zones, as well as the initial and goal positions of the robot, were uniformly sampled within $[-10,10]\times [-10,10]$. To ensure meaningful navigation tasks, the initial and goal positions were sampled at least $7$ m apart and lie outside the static obstacle regions. Each dynamic zone $j \in \mathcal{C}_d$ was assigned a random constant velocity $v_j \in \mathbb{R}^2$ with $||v_j||=2$. The initial conditions were sampled to satisfy all (both hard and soft) constraints.



We compare our methods (Alg.~\ref{alg:iter} and~\ref{alg:local_search}) against two other baselines - Baseline~1 as Chinneck's algorithm~\cite[Algorithm~1]{chinneck} and Baseline~2 as~\cite[Algorithm~2]{hardik2024}. Note that Baseline~2 does not reintroduce the constraints once they are disregarded. However, for a fair comparison in this simulation, we modified it so that the disregarded constraints are allowed to be reintroduced by assigning their corresponding Lagrange multiplier's values from previous step to $0$. Both baseline algorithms evaluate feasibility of a given problem using a slacked linear program (LP). We evaluate performance by comparing computation times and the percentages of disregarded soft constraints across different algorithms in randomly generated environments. The simulations are run on a computer with an Intel(R) Core(TM) i5-10500H CPU @ 2.50GHz and 8.00 GB RAM.


% \begin{table}[h]
% \caption{Comparison between Alg.~\ref{alg:alg1} and baseline algorithms}
% \label{tab:records2}
%     \centering
%     \begin{tabular}{|l|c|c|c|c|c|}
%     \hline & Alg.~\ref{alg:alg1} & \makecell{Baseline~$1$ \\ LP/QP}
% & \makecell{Baseline~$2$ \\ LP/QP} 
%     \\\hline Avg. Comp. Time (s) &  $0.009$ & $0.016$ / $0.051$
%     & $0.016$ / $0.018$
%     \\\hline
%       Max Comp. Time (s) & $0.030$ & $0.126$ / $0.145$ & $0.041$ / $0.045$       \\\hline
%     Avg. Arrival Time (s) & 10.864 & 11.229 & 9.552     \\\hline
%      Max Arrival Time (s) & 12.950 & 14.900 & 11.750 
%       \\\hline Avg. Disregarded (\%) &  $4.85$ & $3.03$ & $9.46$ \\\hline Max Disregarded (\%) & $46.00$ & $24.00$ & $52.00$ \\\hline
% \end{tabular}
% \end{table}

\subsection{Varying Number of Obstacles}

We first performed a set of simulations with a varying number of the number of zones (soft constraints) and a fixed set of $5$ hard constraints corresponding to the CLF and input bounds. The number of zones ranged from $2$ to $100$ in increments of $2$, equally divided between dynamic and static constraints, across $50$ randomly generated environments. We set the search depth $D=5$ for Alg.~\ref{alg:local_search}. 

The results are shown in Fig.~\ref{fig:varying_num_obs}. In general, both of our proposed algorithms achieve performance comparable to Baseline~1 while significantly outperforming Baseline~2 in terms of disregarding constraints. However, our algorithms are in average slightly slower than Baseline~1. Algorithm~\ref{alg:iter} is faster than Baseline~$2$, while Algorithm~\ref{alg:local_search} exhibits similar average computations times to those of Baseline~2. Importantly, both proposed algorithms consistently maintain low computation times across all problem sizes, remaining below $0.07$ seconds even in the largest scenarios. In contrast, baseline algorithms, especially Baseline~1, experience frequent spikes in computation time as shown in Fig.~\ref{fig:varying_num_obs}(b), which may limit their reliabilities for real-time deployment.


\subsection{Fixed Number of Obstacles}

To further examine the performance of the proposed algorithms, we separately conducted simulations with a fixed number of $100$ obstacles across $50$ randomly generated environments. For a more comprehensive evaluation, we varied the search depth $D$ in Alg.~\ref{alg:local_search} among $1$, $5$, and $10$, and included Baseline~1 and Baseline~2 algorithms with quadratic program (QP) feasibility checks. The distributions of disregarded-constraint percentages at each time $t$ for the different algorithms are shown in Fig.~\ref{fig:distribution}. The results are also summarized in Table~\ref{tab:records2}, showcasing the average and maximum computation times, as well as the average and maximum constraint-disregarding rates, for all considered algorithms. 


 These results indicate that Alg.~\ref{alg:iter} and Alg.~\ref{alg:local_search} with different search depths achieve lower constraint-disregarding rates and thus better performance at the cost of only a slight increase in computation time compared to the baseline algorithms with LP feasibility checks, while significantly outperforming the baselines that rely on QP feasibility checks. Furthermore, similar to the previous set of simulations, our algorithms have lower maximum computation times compared to the baseline algorithms.

In summary, the simulation results demonstrate that both of proposed algorithms (Alg.~\ref{alg:iter} and~\ref{alg:local_search}) tend to achieve comparable or even better performance against the baseline algorithms in terms of the number of disregarded constraints. Compared to Baseline~1, the proposed algorithms require a slight increase in computation time, but exhibit greater reliability by consistently maintaining a significantly lower maximum computation times across a wide range of problem sizes. Moreover, unlike the heuristic baseline approaches, our methods are grounded in Cor.~\ref{col:adding_cons} and~\ref{cor:cone_boundary:2}, which provide a theoretically-driven continuous metric quantifying the impact of each constraint on feasibility/infeasibility of a problem. Finally, the proposed formulation solves LPs with only $m$ decision variables, instead of the $n_s + m$ variables required by baseline algorithms with slack variables. This reduction enables more graceful scaling with the number of soft constraints and explains the lower maximum computation times observed in our results.



\begin{figure}
    \centering
    \includegraphics[trim={1.0cm 1.75cm 1.25cm 0.0cm},clip,width=1\linewidth]{figures/fig_2.eps}
    \caption{ Average (a) and Maximum (b) percentages of disregarded constraints, Average (c) and Maximum (d) computation times for Baseline~1, Baseline~2, Alg.~\ref{alg:iter}, and Alg.~\ref{alg:local_search} at each time step~$t$.}
    \label{fig:varying_num_obs}
\end{figure}
\begin{figure}
    \centering
    \includegraphics[trim={1.0cm 0.9cm 1.5cm 1.cm},clip,width=\linewidth]{figures/fig_3.eps}
    \caption{Histograms showing the distribution of the percentage of soft constraints disregarded (out of $50$) at each time instance across $50$ simulations for each algorithm. Frequency (in log-scale) indicates the total number of time instances at which the corresponding percentages of constraints are disregarded across $50$ different simulations.}
    \label{fig:distribution}
\end{figure}

\begin{table}[h]
\caption{Comparison with baseline algorithms}
\label{tab:records2}
    \centering
    \begin{tabular}{|Sl||Sc|Sc|Sc|Sc|Sc|}
    \hline & \makecell{Avg. Comp. \\Time (s)} & \makecell{Max Comp. \\ Time (s)} &  \makecell{Avg.  \\ Drop (\%)} &  \makecell{Max \\ Drop (\%)} \\\hline\hline 
    \makecell{Baseline~$1$ \\ (LP/QP)} &  0.036 /  0.598 & 0.175 / 1.372&  2.245 &     29.000 \\\hline
    \makecell{Baseline~$2$ \\ (LP/QP)} & 0.048 / 0.286
 &   0.091 / 0.403 &  13.580 & 51.000 \\\hline\hline
    \makecell{Alg.~\ref{alg:iter}} & 0.038  &  0.050 &     2.144 & 14.000 \\\hline
   \makecell{Alg.~\ref{alg:local_search} \\ ($D=1$)} & 0.048 & 0.067
 &  2.024 & 15.000 \\\hline
 \makecell{Alg.~\ref{alg:local_search} \\ ($D=5$)} &  0.048    &   0.073 &  2.024 & 15.000 \\\hline
 \makecell{Alg.~\ref{alg:local_search} \\ ($D=10$)} & 0.049  &  0.075
 & 2.024 & 15.000 \\\hline
\end{tabular}
\end{table}


% \begin{algorithm}
% \caption{Local Feasible Configuration Search}\label{alg:feas_perm_srch:2}
% \begin{algorithmic}[1]
% \State 
%  \parbox[t]{\dimexpr\linewidth-\algorithmicindent-20pt}{%
% Given the constraint set \eqref{eq:input_set:config} composed of hard and soft constraints indexed by $\mathcal{C}_h,\mathcal{C}_s$, System \eqref{eq:dynamics:hybrid}, an initial hybrid state $\bar{x} \triangleq  x(t_0,0) = x(0,0)\in\mathbb{R}^n$:}
% \State 
%  \parbox[t]{\dimexpr\linewidth-\algorithmicindent-20pt}{% 
% Set $P(t_0,0) = P(0,0)  = \overline{P}$ according to the hard constraints:}
% \begin{equation}
%     \overline{P} = \begin{cases}
%          1, & i \in \mathcal{C}_h\\
%          0, & i \in \mathcal{C}_s
%     \end{cases}
%     \notag 
% \end{equation}
% \State  Enforce additional constraints through Alg. \ref{alg:iter}:
%  \begin{equation}
%      P{(0,0)} \gets \textrm{\fontfamily{qcr}\selectfont ICA} (A(\bar{x},0),B(\bar{x},0),P{(0,0)}),
%      \notag 
%  \end{equation}
%  \State  Set $j\gets 0$,
% \While{System has not converged}
%     \State 
%      \parbox[t]{\dimexpr\linewidth-\algorithmicindent-20pt}{%
%     Evolve System \eqref{eq:dynamics:hybrid} from the state $x(t_j,j)$ and time $t_j$, under the input $\pi_{\left( A,B\right)}(P,x,t),$}
%     \If{ Prob. \eqref{eq:polar_cone_B:star_time} is infeasible for $P(t,j)$ }
%         \State   Get last feasible solution to \eqref{eq:polar_cone_B:star_time}:
%             \begin{equation}
%                 \nu^\star_{P(t_j,j)}\left(x\left(t_{j+1},j\right),t_j\right),\notag 
%             \end{equation}
%         \State   
%         \parbox[t]{\dimexpr\linewidth-\algorithmicindent-20pt}{%
%         Get the elements of $\nu^\star_{P(t_j,j)}\left(x\left(t_{j+1},j\right),t_j\right)$ corresponding to the enforced constraint $ \nu_e \in \mathbb{R}_{\geq 0}^{n_e}$  (see Thm. \ref{thm:conf_eval}), and sort them, resulting in the index vector $I^{\textrm{e,sorted}}, I^{\textrm{e,sorted}}_k\in \mathcal{C}_s,\forall k\in\{1,\cdots,n_e\}$, where $n_e$ denotes the number of enforced soft constraints,}
%         \State  
%         \parbox[t]{\dimexpr\linewidth-\algorithmicindent-20pt}{%
%          Initialize $P' \gets P(t_{j+1},j), k\gets 1$ and disregard constraint until the problem becomes feasible:}
%         \While{Prob. \eqref{eq:polar_cone_B:star_time} is infeasible for $P'$}
%             \State  $P'_{I^{\textrm{e,sorted}}_k} \gets 0,$
%             \State  $ k\gets k + 1,$
%         \EndWhile
%         \State  Update Configuration;
%             \begin{equation}
%                 P(t_{j+1},j+1) \gets P' ,
%                 \notag 
%             \end{equation}
%         \State  Set $j\gets j+1$,
%     \EndIf
% \EndWhile
% \end{algorithmic}
% \end{algorithm}


% \begin{algorithm}
% \caption{Algorithm 2 + 3}\label{alg:feas_perm_srch:1}
% \begin{algorithmic}[1]
% \State
%  \parbox[t]{\dimexpr\linewidth-\algorithmicindent-20pt}{%
% Given the constraint set \eqref{eq:input_set:config} composed of hard and soft constraints indexed by $\mathcal{C}_h,\mathcal{C}_s$, System \eqref{eq:dynamics:hybrid}, an initial hybrid state $\bar{x} \triangleq  x(t_0,0) = x(0,0)\in\mathbb{R}^n$:}
% \State 
%  \parbox[t]{\dimexpr\linewidth-\algorithmicindent-20pt}{%
% Set $P(t_0,0) = P(0,0) = \begin{cases}
%          1, & i \in \mathcal{C}_h\\
%          0, & i \in \mathcal{C}_s
%     \end{cases}$ according to the hard constraints:}

% \State  Enforce additional constraints through Alg. \ref{alg:iter}:\\
%  \begin{equation}
%      P{(0,0)} \gets \textrm{\fontfamily{qcr}\selectfont ICA} (A(\bar{x},0),B(\bar{x},0),P{(0,0)}),
%      \notag 
%  \end{equation}
%  \State  Set $j\gets 0$,
% \While{System has not converged}
%     \State 
%      \parbox[t]{\dimexpr\linewidth-\algorithmicindent-20pt}{%
%     Evolve System \eqref{eq:dynamics:hybrid} from the state $x(t_j,j)$ and time $t_j$, under the input $\pi_{\left( A,B\right)}(P,x,t),$}
%     \If{ Prob. \eqref{eq:polar_cone_B:star_time} is infeasible for $P(t,j)$ }
%         \State  Re-initialize the configuration by performing lines 8-15 to get $P(t_{j+1},j+1)$.
%         \State  Enforce additional constraints through Alg. \ref{alg:iter}:
%          \begin{equation}
%             \begin{gathered}
%                 P{(t_{j+1},j+1)} \gets \\ 
%                  \textrm{\fontfamily{qcr}\selectfont ICA} (A(P(t_j,j),t_j),B(P(t_j,j),t_j),P{(t_{j+1},j+1)}), 
%             \end{gathered}\notag 
%          \end{equation}
%         \State  Set $j\gets j+1$,
%     \EndIf
% \EndWhile
% \end{algorithmic}
% \end{algorithm}


% \subsection{Feasible Space Monitoring}
% In \cite{11151444}, the authors proposed a scheme for feasible constraint monitoring, i.e., finding inputs to a dynamical system \eqref{eq:dynamics}, such that a set of affine constraints remains non-empty. This was achieved by employing the volume of the respective high-dimensional polytope as a CBF. However, using the volume of polytopes presents two main issues --which the authors acknowledge--: 1) computing a polytope's volume is hard and resource-intensive, 2) the resulting volume is non-smooth, yielding discontinuous gradients in the relevant CBF conditions.  
% \par 
% In order to avoid these issues, note that the solution to \eqref{eq:polar_cone_B:star} can be used instead; we have proved how non-negativity of the polar cone components is a necessary and sufficient condition for feasibility (see Thm. \ref{thm:feas_lp_farkas}, Props. \ref{prop:cone}, \ref{prop:lp_cone_feas}). Therefore, we can directly employ these components, i.e., the solution to \eqref{eq:polar_cone_B:star}, and impose their positivity through CBF conditions. However, non-uniqueness of the nullspace basis in \eqref{eq:polar_cone_B:star} may pose issues of continuity/smoothness. We believe that such continuity issues can be bypassed by computing ``minimally deviating'' basis vectors, or by invoking geometric arguments (e.g., by constructing a map from state/time to a Grassmannian of $K$ planes --where $K$ is the dimension of the kernel of $A(x,t)$). Nevertheless, this requires additional research and is considered outside the scope of this paper.   
% \par 
% Consider again the LP \eqref{eq:feas_lp:star:dual}:
% \begin{equation}
%     \begin{gathered}
%         \nu^\star(x,t) = \underset{\nu\in\mathbb{R}{\color{blue}^c}}{\arg\min}
%         \left\{
%             {0}_C^\top \nu  
%         \right\},        \\
%         \textrm{s.t.: } N_A^\top (x,t)\nu= N^\top_A(x,t)B(x,t), \nu \geq 0,
%     \end{gathered}
%     \tag{\ref{eq:feas_lp:star:dual}}
% \end{equation}
% where $N_A : \mathbb{R}^n\times \mathbb{R}_{\geq 0}\rightarrow{\color{blue}\mathbb{R}^{c\times k}}$ denotes a nullspace basis of the matrix $A(x,t)$ with an overloading of notation.



\section{Conclusion}\label{sec:concl}
In this paper, a theoretical analysis for feasibility assessment of affine constraints was first presented and then employed to formulate online feasible constraint selection algorithms. The proposed framework is demonstrated to provide similar results to state-of-the art competing approaches, while exhibiting reduced computation times. The proposed analysis further paves the way for a number of problems which we intend to pursue in the future, namely investigating solutions to the maxFS problem, creating feasibility-aware/informed controllers and exploring adaptive techniques to render infeasible constrained control problems feasible. These problems have a variety of possible applications for constrained control.   


\bibliographystyle{IEEEtran}
\bibliography{main}
\end{document}


